\documentclass[aps,preprint]{revtex4}%
\usepackage{amsfonts}
\usepackage{amsmath}
\usepackage{amssymb}
\usepackage{graphicx}%
\setcounter{MaxMatrixCols}{30}
%TCIDATA{OutputFilter=latex2.dll}
%TCIDATA{Version=4.10.0.2363}
%TCIDATA{CSTFile=revtex4.cst}
%TCIDATA{Created=Monday, September 20, 2004 09:58:21}
%TCIDATA{LastRevised=Monday, July 25, 2005 11:36:42}
%TCIDATA{<META NAME="GraphicsSave" CONTENT="32">}
%TCIDATA{<META NAME="DocumentShell" CONTENT="Articles\SW\REVTeX 4">}
%TCIDATA{Language=American English}
\newtheorem{theorem}{Theorem}
\newtheorem{acknowledgement}[theorem]{Acknowledgement}
\newtheorem{algorithm}[theorem]{Algorithm}
\newtheorem{axiom}[theorem]{Axiom}
\newtheorem{claim}[theorem]{Claim}
\newtheorem{conclusion}[theorem]{Conclusion}
\newtheorem{condition}[theorem]{Condition}
\newtheorem{conjecture}[theorem]{Conjecture}
\newtheorem{corollary}[theorem]{Corollary}
\newtheorem{criterion}[theorem]{Criterion}
\newtheorem{definition}[theorem]{Definition}
\newtheorem{example}[theorem]{Example}
\newtheorem{exercise}[theorem]{Exercise}
\newtheorem{lemma}[theorem]{Lemma}
\newtheorem{notation}[theorem]{Notation}
\newtheorem{problem}[theorem]{Problem}
\newtheorem{proposition}[theorem]{Proposition}
\newtheorem{remark}[theorem]{Remark}
\newtheorem{solution}[theorem]{Solution}
\newtheorem{summary}[theorem]{Summary}
\newenvironment{proof}[1][Proof]{\noindent\textbf{#1.} }{\ \rule{0.5em}{0.5em}}
\begin{document}
\preprint{ }
\title[Janaks Theorem]{Generalized Janaks Theorem}
\author{B. Weiner}
\affiliation{Pennsylvania State University}
\keywords{}
\pacs{}

\begin{abstract}
.

\end{abstract}
\date{07/18/2005}
\startpage{1}
\maketitle
\tableofcontents


\section{Introduction\label{Intro}}

Janaks theorem is a useful auxiliary tool in the Kuhn-Sham (KS) approach to
Density Functional Theory, that is helpful in analyzing the structure and
meaning of the eigenvalues of the KS Fock operator. It states that these
eigenvalues are equal to the derivatives of the exact total energy with
respect to the occupation numbers of the First Order Reduced Density Operator
(FORDO) of the model KS state. However many conceptual questions surround this
theorem. In what sense can the occupation numbers vary as the model KS state
is by definition an Independent Particle State (IPS)\ with integer occupation
numbers? What is the functional relationship between the eigenvalues and the
occupation number? How does one define a KS\ functional that is differentiable
with respect to occupation numbers?

In both KS-DFT and Hartree-Fock theory these eigenvalues correspond to
Lagrange parameters associated with the normalization constraint on the
orbitals, that constitute the Independent Particle State (IPS) and to
ionization energies of electrons described by these orbitals. This
correspondence can be generalized to Multi Configurational Self Consistent
States (MCSCF), as described for instance in \cite{WeinerOrtiz05}. Constrained
optimization theory \cite{Hestenes} shows that the values of the Lagrange
parameters correspond to the sensitivity of the object function, in this case
the $N$-electron energy, to variations of the constraint, in this case the
normalization of the orbitals. Hence Janaks theorem implies that the
normalization of the orbitals must be a function of the occupation numbers and
hence it would useful to know the form of this relationship.

In order to answer these questions we will review the basic principles of
Hohenburg-Kohn (HK), KS DFT (section 2) and Janaks Theorem (section 3) in the
context of a constrained variational approach and sensitivity theory. However
the Levy-Leib constrained optimization definition of the density functionals
does not lead to an explicit form for these functionals and thus makes their
approximation and investigation difficult.

A well defined KS\ functional determines the total energy in terms of a set of
occupations numbers and one particle orbitals. Such a functional has been
explicitly constructed \cite{Weiner&Ortiz2005} for Antisymmetrized Geminal
Power states in terms of General Spin Orbitals (GSO's) and parameters that
determine the FORDO occupation numbers (see section 4). In section 5 we
investigate the form of the derivatives of this functional with respect to
occupation numbers and the special case of its value at integer occupation. In
section 6 we conclude with a discussion and summary.

\section{Density Functionals \label{Standard}}

In this section we review some fundamental background needed to defined HK and
KS density functionals. We will mention important technical mathematical
aspects that need to be examined, but defer a full analysis of these issues to
another article \cite{weiner05}. We will focus for simplicity on the
mathematically most well behaved cases.

\subsection{Hohenberg-Kohn DFT}

The HK density functional is defined in terms of the kinetic energy operator,
$K,$ the coulomb operator, $V_{2},$ and the charge density $\rho$.

\subsubsection{Densities and States}

\bigskip Density Functional Theory is based on the existence of a universal
functional, $\mathcal{\tilde{F}}\left(  \rho\right)  $ \cite{HK} of the charge
density $\rho\left(  \mathbf{r}\right)  ,$ where
\begin{equation}
\rho\left(  \mathbf{r}\right)  =\int D^{1}\left(  \mathbf{r,\xi;r,\xi}\right)
d\mathbf{\xi}%
\end{equation}
and $D^{1}\left(  \mathbf{r,\xi;r}^{\prime}\mathbf{,\xi}^{\prime}\right)  $ is
the integral kernel of the First Order Reduced Operator (FORDO) $D^{1}$
\cite{Coleman} (see Appendix \ref{DenMat} for details of density and reduced
density operators and matrices), which can be expanded as
\begin{equation}
D^{1}=\int D^{1}\left(  \mathbf{r,\xi;r}^{\prime}\mathbf{,\xi}^{\prime
}\right)  \left\vert \mathbf{r,\xi}\right\rangle \left\langle \mathbf{r}%
^{\prime}\mathbf{,\xi}^{\prime}\right\vert d\mathbf{r}d\mathbf{r}^{\prime
}d\mathbf{\xi}d\mathbf{\xi}^{\prime}%
\end{equation}
in the basis $\left\{  \left\vert \mathbf{r,\xi}\right\rangle \right\}  $
formed by the tensor product of the generalized eigenvectors, $\left\{
\left\vert \mathbf{r}\right\rangle \right\}  ,$ of the position operator and a
family of $SU(2)$-coherent spin states $\left\{  \left\vert \mathbf{\xi
}\right\rangle \right\}  $ $\cite{weiner2005}$ . In order to consider spin
dependent effects it is convenient to extend the original formulation of DFT
in terms of $\rho\left(  \mathbf{r}\right)  $ and consider a universal
functional, $\mathcal{F}\left(  \gamma\right)  $, of the space-spin density
$\gamma\left(  \mathbf{r,\xi}\right)  =$ $D^{1}\left(  \mathbf{r,\xi;r,\xi
}\right)  $ \cite{weiner-trickey} and local external potentials of the
position variable $\mathbf{r}$ and spin space variable $\mathbf{\xi}$ that
describe the interaction with external electric and magnetic fields. $\lceil
$The integral kernel of the magnetic dipole operator is local (i.e. diagonal)
in the coherent state representation but not over the spectrum of the spin
operator$.\rfloor$

For a system that contains on average $N$-electrons the space spin density
must satisfy the condition
\begin{equation}
\int\gamma\left(  \mathbf{r,\xi}\right)  d\mathbf{r}d\mathbf{\xi}=N,
\label{Number}%
\end{equation}
which does not restrict the density to be the contraction of a pure
$N$-electron state and one could consider incoherent and coherent mixtures of
sates describing different number of electrons. Focusing on pure $N$-electron
states one can define convex sets, $\left\{  \left[  \gamma\right]
_{N}\right\}  ,$ of $N$-electron state density operators by
\begin{equation}
\left[  \gamma\right]  _{N}=\left\{  \left.  D\right\vert \Xi_{1}\left(
D\right)  =\gamma;D\in S_{N}\subset\mathcal{B}_{1}\left(  \mathcal{H}%
^{N}\right)  \right\}
\end{equation}
(see Appendix \ref{DenMat} for the definition of the domain $S_{N}$ and
Appendix \ref{Operators} for the definition of $\mathcal{B}_{1}\left(
\mathcal{H}^{N}\right)  $). The constraining functional $\Xi_{1}$ can be
written in terms of the set of density operators $\left\{  \left.  \Phi\left(
\mathbf{x}\right)  ^{\dagger}\Phi\left(  \mathbf{x}\right)  \right\vert
\mathbf{x\in}\mathbb{X}\right\}  $ as
\begin{equation}
\Xi_{1}\left(  D\right)  \left(  \mathbf{x}\right)  =\operatorname{Tr}\left\{
\Phi\left(  \mathbf{x}\right)  ^{\dagger}\Phi\left(  \mathbf{x}\right)
D\right\}  =\gamma\left(  \mathbf{x}\right)  .
\end{equation}
One can view this set of density operators as a generalized vector operator
$\Phi^{\dagger}\Phi\mathbf{\ }$with a continuum of components. The
satisfaction of the condition eq. $\left(  \ref{Number}\right)  $ for
$D\in\mathcal{B}_{1}\left(  \mathcal{H}^{N}\right)  $ implies that the state
is normalized (see Appendix \ref{Normalization}) i.e.
\begin{equation}
\operatorname{Tr}D=1,
\end{equation}
and hence $D$ needs only to be constrained to lie in the cone (see Appendix
\ref{Operators}) of positive operators of $\mathcal{B}_{1}\left(
\mathcal{H}^{N}\right)  $. As $D$ is a positive operator it can always be
expressed as the product
\begin{equation}
D=QQ^{\dagger},
\end{equation}
where $\left(  \text{Appendix \ref{Operators}}\right)  $
\begin{equation}
Q\in\mathcal{B}_{2}\left(  \mathcal{H}^{N}\right)  ,
\end{equation}
and $\mathcal{B}_{2}\left(  \mathcal{H}^{N}\right)  $ is the Hilbert space of
Hilbert-Schmidt operators acting on $N$-electron space $\mathcal{H}^{N}$.
However this decomposition is not unique, as in particular
\begin{equation}
D=QUU^{\dagger}Q^{\dagger}%
\end{equation}
for any unitary operator $U\in\mathcal{U}\left(  \mathcal{H}^{N}\right)  $.
One can remove this unitary group redundancy by choosing the operator $Q$ to
be self adjoint, leading to
\begin{equation}
D=A^{2},
\end{equation}
where
\begin{equation}
A\in\mathcal{B}_{2}\left(  \mathcal{H}^{N}\right)  \text{ and }A=A^{\dagger}.
\end{equation}
The operator $A$ is the square root of $D$, which is not unique, as can be
observed from spectral expansion
\begin{equation}
A=D^{\frac{1}{2}}=\sum_{1\leq j\leq r_{N}}\pm\sqrt{\alpha_{j}}\left\vert
\Upsilon_{j}\right\rangle \left\langle \Upsilon_{j}\right\vert , \label{specA}%
\end{equation}
where
\begin{equation}
D=\sum_{1\leq j\leq r_{N}}\alpha_{j}\left\vert \Upsilon_{j}\right\rangle
\left\langle \Upsilon_{j}\right\vert ;\,\alpha_{j}\geq0,1\leq j\leq r_{N}%
\end{equation}
and $\left\{  \left.  \left\vert \Upsilon_{j}\right\rangle \right\vert 1\leq
j\leq r_{N}\right\}  $ is a complete orthornormal set of eigenvectors of $D$.
One could restrict $A$ to be positive i.e. use the positive square roots in
eq. $\left(  \ref{specA}\right)  ,$ but that would be too restrictive for our
purposes, as the set of positive Hilbert Schmidt operators forms a cone but do
not form a linear space. However, the multi valuedness of the square root
function is not a problem as the different roots are separated from each other
and could only form isolated local minima. The domain of variation of the
$N$-electron energy function can thus be taken to be the \emph{real} Hilbert
space $\mathcal{B}_{2sa}\left(  \mathcal{H}^{N}\right)  ,$ of dimension
$r_{N}^{2},$ which has an orthonormal basis
\begin{align}
\Lambda_{jk}  &  =\frac{1}{\sqrt{2}}\left\{  \left\vert \Psi_{j}\right\rangle
\left\langle \Psi_{k}\right\vert +\left\vert \Psi_{k}\right\rangle
\left\langle \Psi_{j}\right\vert \right\}  ;\quad1\leq j<k\leq r_{N}%
\nonumber\\
\Lambda_{jj}  &  =\left\vert \Psi_{j}\right\rangle \left\langle \Psi
_{j}\right\vert ;\quad1\leq j\leq r_{N}\nonumber\\
\Lambda_{jk}  &  =\frac{i}{\sqrt{2}}\left\{  \left\vert \Psi_{j}\right\rangle
\left\langle \Psi_{k}\right\vert -\left\vert \Psi_{k}\right\rangle
\left\langle \Psi_{j}\right\vert \right\}  ;\quad1\leq j>k\leq r_{N},
\end{align}
where $\left\{  \left.  \left\vert \Psi_{j}\right\rangle \right\vert 1\leq
j\leq r_{N}\right\}  $ is any conb for $\mathcal{H}^{N},$ of course all these
spaces are infinite dimensional if $r_{N}=\infty.$

\subsubsection{The Density Constraint}

The Levy-Leib construction is based on defining the value of the density
functional as the minimum of a energy variation over states constrained to
have a fixed density. The density constraint can either be formulated in terms
of a quadratic functional $\Xi_{2}$ or a linear functional $\Xi_{1}.$ In terms
of the quadratic functional
\begin{equation}
\Xi_{2}:\mathcal{B}_{2sa}\left(  \mathcal{H}^{N}\right)  \rightarrow
L_{1+}\left(  \mathbb{X}\right)  \subset L_{1}\left(  \mathbb{X}\right)  ,
\end{equation}
where $L_{1+}\left(  \mathbb{X}\right)  $ is the positive cone of
$L_{1}\left(  \mathbb{X}\right)  ,$ it is given by
\begin{equation}
\Xi_{2}\left(  A\right)  =\operatorname{Tr}\left\{  A\Phi^{\dagger}\Phi
A\right\}  =\left(  A\left\vert \Phi^{\dagger}\Phi A\right.  \right)  =\gamma,
\end{equation}
where the real symmetric inner product $\left(  \cdot\left\vert \cdot\right.
\right)  $ in $\mathcal{B}_{2sa}\left(  \mathcal{H}^{N}\right)  $ is defined
as
\begin{equation}
\left(  X\left\vert Y\right.  \right)  =\operatorname{Tr}\left\{  XY\right\}
.
\end{equation}
In terms of the linear functional $\Xi_{1}$%
\begin{equation}
\Xi_{1}:\mathcal{B}_{1sa}\left(  \mathcal{H}^{N}\right)  \rightarrow
L_{1}\left(  \mathbb{X}\right)
\end{equation}
it is given by
\begin{equation}
\Xi_{1}\left(  X\right)  =\operatorname{Tr}\left\{  \Phi^{\dagger}\Phi
X\right\}  =\gamma,
\end{equation}
where
\begin{equation}
X\geq0.
\end{equation}
If the differentials of the components of the constraints at a point $X$ form
a basis (see Appendix basis \ref{Basis}) for $L_{1}\left(  \mathbb{X}\right)
$ then this point is called \emph{regular}. The components are
\begin{align}
D\Xi_{2}\left(  X\right)  \left(  \mathbf{x}\right)   &  =\operatorname{Tr}%
\left\{  X\Phi\left(  \mathbf{x}\right)  ^{\dagger}\Phi\left(  \mathbf{x}%
\right)  \cdot\right\}  =\left(  X\left\vert \Phi\left(  \mathbf{x}\right)
^{\dagger}\Phi\left(  \mathbf{x}\right)  \right.  \cdot\right) \\
D\Xi_{1}\left(  X\right)  \left(  \mathbf{x}\right)   &  =\Xi_{1}\left(
\mathbf{x}\right)  =\operatorname{Tr}\left\{  \Phi\left(  \mathbf{x}\right)
^{\dagger}\Phi\left(  \mathbf{x}\right)  \cdot\right\}  .
\end{align}


It will be convenient to use multiplication super operators
\begin{equation}
\widehat{T}:\mathcal{B}_{2sa}\left(  \mathcal{H}^{N}\right)  \rightarrow
\mathcal{B}_{2sa}\left(  \mathcal{H}^{N}\right)  ,
\end{equation}
which could be unbounded, and are defined by the quadratic form
\cite{weiner05}
\begin{equation}
\left(  X\left\vert \widehat{T}Y\right.  \right)  =\operatorname{Tr}\left\{
XTY\right\}  \,\forall X\in\mathcal{B}_{2sa}\left(  \mathcal{H}^{N}\right)
,\,Y\in\operatorname{Dom}\left\{  T\right\}  ,
\end{equation}
giving
\begin{equation}
\widehat{T}\left\vert Y\right)  =\left\vert \left[  TY\right]  \right)  ,
\end{equation}
where
\begin{equation}
T:\operatorname{Dom}\left\{  T\right\}  \subseteq\mathcal{H}^{N}%
\rightarrow\mathcal{H}^{N}%
\end{equation}
and $\left\vert \left[  TY\right]  \right)  $ is the projection of $TY$ onto
$\mathcal{B}_{2sa}\left(  \mathcal{H}^{N}\right)  $ i.e. the self adjoint part
of $TY,$ given by
\[
\left[  TY\right]  =\frac{1}{2}\left\{  TY+YT\right\}  =\frac{1}{2}\left[
T,Y\right]  _{+}.
\]
One can observe that
\begin{equation}
\operatorname{Tr}\left\{  XTY\right\}  =\operatorname{Tr}\left\{  X\left[
TY\right]  \right\}  ,
\end{equation}
as $X\in\mathcal{B}_{2sa}\left(  \mathcal{H}^{N}\right)  $ and is orthogonal
to the subspace of anti self adjoint operators.

\subsubsection{Hohenberg-Kohn Density Functional Construction}

Even though the explicit form of the Hohenberg-Kohn functional $\mathcal{F}%
_{HK}\left(  \gamma\right)  $ is not known one can define it implicitly
\cite{LevyLieb} as
\begin{equation}
\mathcal{F}_{HK}\left(  \gamma\right)  =\operatorname{Tr}\left\{  \left(
K+V_{2}\right)  D_{\ast}\left(  \gamma\right)  \right\}  =E_{U}\left(
D_{\ast}\left(  \gamma\right)  \right)
\end{equation}
in terms of the minima, $E_{U}\left(  D_{\ast}\left(  \gamma\right)  \right)
,$ of the universal energy functional $E_{U}\left(  D\right)  ,$ that are
solutions of the constrained linear optimization
\begin{equation}
\min_{D\in\mathcal{B}_{1sa}\left(  \mathcal{H}^{N}\right)  }\operatorname{Tr}%
\left\{  \left(  K+V_{2}\right)  D\right\}  \equiv\min_{D\in\mathcal{B}%
_{1sa}\left(  \mathcal{H}^{N}\right)  }\text{ }E_{U}\left(  D\right)
\label{HK}%
\end{equation}
subject to
\begin{align}
\Xi_{1}\left(  D\right)   &  =\gamma\nonumber\\
D  &  \geq0, \label{HKC}%
\end{align}
which is a linear programing problem, or by the constrained quadratic
optimization
\begin{equation}
\min_{A\in\mathcal{B}_{2sa}\left(  \mathcal{H}^{N}\right)  }\operatorname{Tr}%
\left\{  A\overset{}{\left(  K+V_{2}\right)  A}\right\}  =\left(  A_{\ast
}\left(  \gamma\right)  \left\vert \widehat{\left(  K+V_{2}\right)  }A_{\ast
}\left(  \gamma\right)  \right.  \right)  =E_{U}\left(  A_{\ast}\left(
\gamma\right)  \right)  \label{HKQ}%
\end{equation}
subject to
\begin{equation}
\Xi_{2}\left(  A\right)  =\left(  A\left\vert \widehat{\Phi^{\dagger}\Phi
}A\right.  \right)  =\gamma, \label{HKQC}%
\end{equation}
where
\begin{equation}
D_{\ast}\left(  \gamma\right)  =A_{\ast}\left(  \gamma\right)  ^{2}.
\end{equation}
We will focus on the quadratic formulation, which is a convex constrained
optimization problem. The Lagrangian associated with this constrained
optimization problem eq. $\left(  \ref{HK}\right)  $ is
\begin{equation}
\mathcal{L}\left(  A,\lambda,\gamma\right)  =E_{U}\left(  A\right)
-\left\langle \left.  \lambda\right\vert \left\{  \Xi_{2}\left(  A\right)
-\gamma\right\}  \right\rangle , \label{Lar}%
\end{equation}
where the universal energy is
\begin{equation}
E_{U}\left(  A\right)  =\operatorname{Tr}\left\{  \left(  K+V_{2}\right)
A^{2}\right\}  =\left(  A\left\vert \widehat{\left(  K+V_{2}\right)
}A\right.  \right)  ,
\end{equation}
the Lagrange multiplier function is
\begin{equation}
\left\langle \left.  \lambda\right\vert \left\{  \Xi_{2}\left(  A\right)
-\gamma\right\}  \right\rangle =\int\lambda\left(  \mathbf{x}\right)  \left\{
\Xi_{2}\left(  A\right)  \left(  \mathbf{x}\right)  -\gamma\left(
\mathbf{x}\right)  \right\}  d\mathbf{x,\,}\lambda\in L_{\infty}\left(
\mathbb{X}\right)  ,\gamma\in L_{1}\left(  \mathbb{X}\right)
\end{equation}
where $\left\langle \cdot|\cdot\right\rangle $ is the scalar product between
the dual spaces $\left\langle L_{\infty}\left(  \mathbb{X}\right)
,L_{1}\left(  \mathbb{X}\right)  \right\rangle $ i.e.
\begin{equation}
\left\langle \cdot|\cdot\right\rangle :L_{\infty}\left(  \mathbb{X}\right)
\times L_{1}\left(  \mathbb{X}\right)  \rightarrow\mathbb{R}%
\end{equation}
and the control density $\gamma$ is chosen to belong to the convex set
$\mathcal{P}_{N}$ given by
\begin{equation}
\mathcal{P}_{N}=\left\{  f\left\vert f\in L_{1}\left(  \mathbb{X}\right)
;\,f\left(  \mathbf{x}\right)  \geq0;\int f\left(  \mathbf{x}\right)  d\left(
\mathbf{x}\right)  =N\right.  \right\}  .
\end{equation}
This choice of the density $\gamma$ coupled with insisting that $A\in
\mathcal{B}_{2sa}\left(  \mathcal{H}^{N}\right)  $ ensures that $A$ is
normalized i.e.
\begin{equation}
\left(  A\left\vert A\right.  \right)  =1.
\end{equation}


The constraint term in the Lagrangian eq. $\left(  \ref{Lar}\right)  $ can be
developed to give
\begin{equation}
\left\langle \lambda\left\vert \gamma\right.  \right\rangle -\left(
A\left\vert \widehat{\Omega\left(  \lambda\right)  }\right.  A\right)
=\left(  A\left\vert \left\{  \left\langle \lambda\left\vert \gamma\right.
\right\rangle \hat{I}-\widehat{\Omega\left(  \lambda\right)  }\right\}
\right.  A\right)  ,
\end{equation}
where $\Omega\left(  \lambda\right)  $ is a local potential defined as
\begin{equation}
\Omega\left(  \lambda\right)  =\int\lambda\left(  \mathbf{x}\right)
\mathbf{\Phi}\left(  \mathbf{x}\right)  ^{\dagger}\mathbf{\Phi}\left(
\mathbf{x}\right)  d\mathbf{x}%
\end{equation}
and
\begin{equation}
\left(  A\left\vert \widehat{\Omega\left(  \lambda\right)  }\right.  A\right)
=\int\lambda\left(  \mathbf{x}\right)  \Xi_{2}\left(  A\right)  \left(
\mathbf{x}\right)  d\mathbf{x.}%
\end{equation}


\subsubsection{Optimization Conditions}

Regular solutions $A_{\ast}$ of the constrained quadratic optimization eqs.
$\left(  \ref{HKQ}\right)  $ and $\left(  \ref{HKQC}\right)  $ correspond to
unconstrained local saddle points $\left(  A_{\ast},\lambda_{\ast}\right)  $
\cite{KKT} of the Lagrangian $\mathcal{L}\left(  A_{\ast},\lambda_{\ast
},\gamma\right)  $
\begin{equation}
\mathcal{L}\left(  A_{\ast},\lambda_{\ast},\gamma\right)  =\max_{\lambda}%
\min_{A}\mathcal{L}\left(  A,\lambda,\gamma\right)
\end{equation}
that satisfy the necessary condition
\begin{equation}
D\mathcal{L}\left(  A_{\ast},\lambda_{\ast},\gamma\right)  =0,
\end{equation}
where the linear map
\begin{equation}
D\mathcal{L}\left(  A_{\ast},\lambda_{\ast},\gamma\right)  :\mathcal{B}%
_{2sa}\left(  \mathcal{H}^{N}\right)  \times L_{\infty}\left(  \mathbb{X}%
\right)  \rightarrow\mathbb{R}%
\end{equation}
is given by%
\begin{equation}
D\mathcal{L}\left(  A_{\ast},\lambda_{\ast},\gamma\right)  \left(
h,\varphi\right)  =\left\langle \varphi\left\vert \gamma\right.  \right\rangle
-\left(  A_{\ast}\left\vert \widehat{\Omega\left(  \varphi\right)  }\right.
A_{\ast}\right)  +2\left(  h\left\vert \left\{  \widehat{K+V_{2}-\Omega\left(
\lambda\right)  }\right\}  \right.  A_{\ast}\right)  . \label{LagDiff}%
\end{equation}
As $A_{\ast}$ is a regular point of the constraint $\lambda_{\ast}$ is unique.
The differential eq. $\left(  \ref{LagDiff}\right)  $ can be partitioned into
$\left(  D_{AA}\mathcal{L}\left(  A,\lambda,\gamma\right)  ,D_{\lambda
}\mathcal{L}\left(  A,\lambda,\gamma\right)  \right)  $, where
\begin{align}
D_{A}\mathcal{L}\left(  A,\lambda,\gamma\right)   &  =2\left(  \cdot\left\vert
\widehat{\left\{  K+V_{2}-\Omega\left(  \lambda\right)  \right\}  }\right.
A\right) \nonumber\\
D_{\lambda}\mathcal{L}\left(  A,\lambda,\gamma\right)   &  =\left\langle
\cdot\left\vert \gamma\right.  \right\rangle -\left(  A\left\vert
\Omega\left(  \cdot\right)  \right.  A\right)  ,
\end{align}
thus giving the stationary point conditions
\begin{subequations}
\begin{align}
\left(  h\left\vert \widehat{\left\{  K+V_{2}-\Omega\left(  \lambda_{\ast
}\right)  \right\}  }\right.  A_{\ast}\right)   &  =0\,\forall h\in
\mathcal{B}_{2sa}\left(  \mathcal{H}^{N}\right) \label{stata}\\
\left\langle \varphi\left\vert \gamma\right.  \right\rangle -\left(  A_{\ast
}\left\vert \widehat{\Omega\left(  \varphi\right)  }\right.  A_{\ast}\right)
&  =0\,\forall\varphi\in L_{\infty}\left(  \mathbb{X}\right)  . \label{statb}%
\end{align}
If eq. $\left(  \ref{stata}\right)  $ and eq. $\left(  \ref{statb}\right)  $
are true for all operators $h$ belonging to $\mathcal{B}_{2sa}\left(
\mathcal{H}^{N}\right)  $ and all functions $\varphi$ belonging to $L_{\infty
}\left(  \mathbb{X}\right)  $ these conditions are equivalent to
\end{subequations}
\begin{equation}
\widehat{\left\{  K+V_{2}-\Omega\left(  \lambda_{\ast}\right)  \right\}
}\left\vert A_{\ast}\right)  =0 \label{Stateig}%
\end{equation}
and constraining condition
\begin{equation}
\left(  A_{\ast}\left\vert \widehat{\Phi^{\dagger}\Phi}\right.  A_{\ast
}\right)  =\gamma.
\end{equation}
(For unbounded operators $K+V_{2}-\Omega\left(  \lambda_{\ast}\right)  $ one
must be more precise, but we will only consider the case when an eigenvector
$\left\vert A_{\ast}\right)  $ exists and thus eq. $\left(  \ref{stata}%
\right)  $ is true for all $h\in\mathcal{B}_{2sa}\left(  \mathcal{H}%
^{N}\right)  $). The optimal value of the Lagrangian satisfies the following
identities:
\begin{align}
\mathcal{L}\left(  A_{\ast},\lambda_{\ast},\gamma\right)   &  =\left(
A_{\ast}\left(  \gamma\right)  \left\vert \left\{  \widehat{\left\{
K+V_{2}-\Omega\left(  \lambda_{\ast}\right)  \right\}  }+\left\langle \left.
\lambda_{\ast}\left(  \gamma\right)  \right\vert \gamma\right\rangle \hat
{I}\right\}  \right.  A_{\ast}\left(  \gamma\right)  \right)  =\left\langle
\left.  \lambda_{\ast}\left(  \gamma\right)  \right\vert \gamma\right\rangle
\nonumber\\
&  =\left(  A_{\ast}\left(  \gamma\right)  \left\vert \widehat{K+V_{2}%
}\right.  A_{\ast}\left(  \gamma\right)  \right)  =E_{U}\left(  A_{\ast
}\left(  \gamma\right)  \right)  .
\end{align}
This shows that the energy of the universal energy functional at the density
$\gamma$ is equal to the value of the Lagrangian and is also the value of the
HK energy functional i.e.
\begin{equation}
\mathcal{F}_{HK}\left(  \gamma\right)  =E_{U}\left(  A_{\ast}\left(
\gamma\right)  \right)  =\left\langle \left.  \lambda_{\ast}\left(
\gamma\right)  \right\vert \gamma\right\rangle .
\end{equation}


The Hessian $D^{2}\mathcal{L}\left(  A,\lambda,\gamma\right)  $ of the
Lagrangian at the point $\left(  A,\lambda,\gamma\right)  $ is given by
\begin{equation}
D^{2}\mathcal{L}\left(  A,\lambda,\gamma\right)  =\left(
\begin{array}
[c]{cc}%
D_{AA}^{2}\mathcal{L}\left(  A,\lambda,\gamma\right)  & D_{A\lambda}%
^{2}\mathcal{L}\left(  A,\lambda,\gamma\right) \\
D_{\lambda A}^{2}\mathcal{L}\left(  A,\lambda,\gamma\right)  & D_{\lambda}%
^{2}\mathcal{L}\left(  A,\lambda,\gamma\right)
\end{array}
\right)  =\left(
\begin{array}
[c]{cc}%
D_{AA}^{2}\mathcal{L}\left(  A,\lambda,\gamma\right)  & D_{A\lambda}%
^{2}\mathcal{L}\left(  A,\lambda,\gamma\right) \\
D_{\lambda A}^{2}\mathcal{L}\left(  A,\lambda,\gamma\right)  & \mathbb{O}%
\end{array}
\right)  ,
\end{equation}
where%
\begin{align}
D_{AA}^{2}\mathcal{L}\left(  A,\lambda,\gamma\right)   &  :\mathcal{B}%
_{2sa}\left(  \mathcal{H}^{N}\right)  \times\mathcal{B}_{2sa}\left(
\mathcal{H}^{N}\right)  \rightarrow\mathbb{R}\nonumber\\
D_{\lambda A}^{2}\mathcal{L}\left(  A,\lambda,\gamma\right)   &  :L_{\infty
}\left(  \mathbb{X}\right)  \times\mathcal{B}_{2sa}\left(  \mathcal{H}%
^{N}\right)  \rightarrow\mathbb{R}%
\end{align}
and%
\begin{align}
D_{AA}^{2}\mathcal{L}\left(  A,\lambda,\gamma\right)   &  =2\left(
\cdot\left\vert \widehat{\left\{  K+V_{2}-\Omega\left(  \lambda\right)
\right\}  }\right.  \cdot\right) \nonumber\\
D_{A\lambda}^{2}\mathcal{L}\left(  A,\lambda,\gamma\right)   &  =-2\left(
A\widehat{\left\vert \Omega\left(  \cdot\right)  \right.  }\cdot\right)
=-2\left(  A\left\vert \widehat{\Phi^{\dagger}\Phi}\right.  \cdot\right)
=-2D_{A}\Xi_{2}\left(  A\right)  .
\end{align}
A stationary point $A_{\ast}$ is a constrained local minimum of $E_{U}$ if the
sub block $D_{A}^{2}\mathcal{L}\left(  A_{\ast},\lambda_{\ast},\gamma\right)
$ of the Hessian $D^{2}\mathcal{L}\left(  A_{\ast},\lambda_{\ast}%
,\gamma\right)  $ is positive at the point $A$ i.e.%
\begin{equation}
\left(  h\left\vert \widehat{\left\{  K+V_{2}-\Omega\left(  \lambda_{\ast
}\right)  \right\}  }\right.  h\right)  \geq0\,\forall h\in\mathcal{B}%
_{2sa}\left(  \mathcal{H}^{N}\right)  .
\end{equation}
It is non degenerate, i.e. unique, if this is a strict inequality. At a
regular point $A$ of the constraint
\begin{equation}
\Xi_{2}\left(  A\right)  :\mathcal{B}_{2sa}\left(  \mathcal{H}^{N}\right)
\rightarrow L_{1}\left(  \mathbb{X}\right)
\end{equation}
the differential $D\Xi_{2}\left(  A\right)  $ of the functional $\Xi_{2}$
\begin{equation}
D\Xi_{2}\left(  A\right)  :\mathcal{B}_{2sa}\left(  \mathcal{H}^{N}\right)
\rightarrow L_{1}\left(  \mathbb{X}\right)
\end{equation}
is surjective \cite{weiner2005}. If the universal energy is differentiable at
an optimal solution $A_{\ast}$ i.e. $DE_{U}\left(  A_{\ast}\right)  $ exists
and if $A_{\ast}$ is regular point of the constraint, then it is a regular solution.

If two states $D_{\ast1}\left(  \gamma\right)  =$ $A_{\ast1}\left(
\gamma\right)  ^{2}$ and $D_{\ast2}\left(  \gamma\right)  =$ $A_{\ast2}\left(
\gamma\right)  ^{2}$ are degenerate then the states
\begin{equation}
\alpha_{1}D_{\ast1}\left(  \gamma\right)  +\alpha_{2}D_{\ast2}\left(
\gamma\right)  =\left(  1-\alpha\right)  D_{\ast1}\left(  \gamma\right)
+\alpha D_{\ast2}\left(  \gamma\right)
\end{equation}
are also degenerate, so one has a line of minima, generated by $\left\{
D_{\ast1}\left(  \gamma\right)  ,D_{\ast2}\left(  \gamma\right)  \right\}  ,$
which are not isolated. Hence%
\begin{equation}
\operatorname{Tr}\left\{  \left\{  K+V_{2}-\Omega\left(  \lambda_{\ast}\left(
\gamma\right)  \right)  \right\}  D_{\ast1}\left(  \gamma\right)  \right\}
=\operatorname{Tr}\left\{  \left\{  K+V_{2}-\Omega\left(  \lambda_{\ast
}\left(  \gamma\right)  \right)  \right\}  D_{\ast2}\left(  \gamma\right)
\right\}  =0.
\end{equation}
As $D_{\ast1}\left(  \gamma\right)  $ and $D_{\ast2}\left(  \gamma\right)  $
are unconstrained stationary states of $K+V_{2}-\Omega\left(  \lambda_{\ast
}\left(  \gamma\right)  \right)  $
\begin{equation}
\left[  \left\{  K+V_{2}-\Omega\left(  \lambda_{\ast}\left(  \gamma\right)
\right)  \right\}  ,D_{\ast j}\left(  \gamma\right)  \right]  =0.
\end{equation}
This implies that
\begin{equation}
\left[  \left\{  K+V_{2}-\Omega\left(  \lambda_{\ast}\left(  \gamma\right)
\right)  \right\}  ,A_{\ast j}\left(  \gamma\right)  \right]  =0.
\end{equation}
As
\begin{equation}
\widehat{\left\{  K+V_{2}-\Omega\left(  \lambda_{\ast}\left(  \gamma\right)
\right)  \right\}  }\left\vert A_{\ast j}\left(  \gamma\right)  \right)
=0\Rightarrow\left[  \left\{  K+V_{2}-\Omega\left(  \lambda_{\ast}\left(
\gamma\right)  \right)  \right\}  ,A_{\ast j}\left(  \gamma\right)  \right]
_{+}=0,
\end{equation}
we can deduce that
\begin{equation}
\left\{  K+V_{2}-\Omega\left(  \lambda_{\ast}\left(  \gamma\right)  \right)
\right\}  A_{\ast j}\left(  \gamma\right)  =0.
\end{equation}
$\lceil$As
\begin{align}
D_{1}  &  =A_{1}^{2}\\
D_{2}  &  =A_{2}^{2}%
\end{align}
and $A_{1}$ and $A_{2}$ are degenerate, but
\begin{equation}
\beta_{1}A_{1}+\beta_{2}A_{2}%
\end{equation}
are not minima unless
\begin{equation}
A_{1}A_{2}=0.\rfloor
\end{equation}


\subsubsection{Mathematical Qualifications}

The existence of the differential of $\left(  A\left\vert \widehat{\left(
K+V_{2}\right)  }A\right.  \right)  $ at a point $A$ is equivalent to Frechet
differentiability at the point $A$. If $\widehat{\left(  K+V_{2}\right)  }$ is
unbounded then this quadratic form is not continuous at all points and hence
there are points where it is not Frechet differentiable and thus where the
differential does not exist. At these points we cannot use the eigenvector
expression eq. $\left(  \ref{Stateig}\right)  .$ However the function $\left(
A\left\vert \widehat{\left(  K+V_{2}\right)  }A\right.  \right)  $ does have a
minimal value $A_{\ast}\left(  \gamma\right)  $ for each $\gamma$ and
\begin{align}
\left(  h\left\vert \widehat{\left(  K+V_{2}\right)  }A_{\ast}\left(
\gamma\right)  \right.  \right)   &  =0\\
\left(  h\left\vert \widehat{\left(  K+V_{2}\right)  }h\right.  \right)   &
\geq0
\end{align}
for all operators $h$ belonging to the domain of $\widehat{\left(
K+V_{2}\right)  }.$ If $ht$ belongs to the domain of $\widehat{\left(
K+V_{2}\right)  }$ for $t\leq\left\vert \delta\right\vert $ and the quadratic
form is continuous at these points then one can produce the Gateux derivative
in the direction of $h.$ At such points a weak form of the Lagrange multiplier
theorem holds.

Another point to consider is when $\widehat{\left(  K+V_{2}\right)  }$ has
only a continuous spectrum, then no eigenvector exists, even though a minimum
does. Again a weak form of the Lagrange multiplier theorem holds.

These considerations have implications about $V$-representatbility \cite{Vrep}

A further phenomena is when the local minimum is degenerate, then the point
$A_{\ast}$ does not produce an invertible Hessian sub block $D_{AA}%
^{2}\mathcal{L}$ so a unique $D_{\ast}$ $=A_{\ast}^{2}$ does not exist. This
does not effect the definition of the HK energy functional, but it does the
definitions of the HK kinetic energy and XC functionals as different
degenerate $D_{\ast}$'s produce different values for these component
functionals. The multivaluedness of the square root function $D_{\ast}%
^{\frac{1}{2}}$ produce isolated local minima that are degenerate.

\subsubsection{Sensitivity Analysis}

Sensitivity analysis \cite{sensitivity} examines how the optimal solution
changes with variation of the value of the constraining functions, which can
be thought of as control paramaters. In this case the $N$-electron state
energy with variations of $\gamma.$ If the optimal points $\left\{  A_{\ast
},\lambda_{\ast},\gamma\right\}  $ are regular for a range of $\gamma^{\prime
}s$ then $\left\{  A_{\ast},\lambda_{\ast}\right\}  $ are in fact functions of
$\gamma$ \cite{weiner2005} i.e.
\begin{align}
A_{\ast}  &  =A_{\ast}\left(  \gamma\right) \nonumber\\
\lambda_{\ast}  &  =\lambda_{\ast}\left(  \gamma\right)  ,
\end{align}
hence
\begin{equation}
D_{\ast}=D_{\ast}\left(  \gamma\right)  =A_{\ast}\left(  \gamma\right)  ^{2}.
\end{equation}
As all points along these paths are feasable one obtains the important
identity
\begin{equation}
\mathcal{L}\left(  A_{\ast},\lambda_{\ast},\gamma\right)  =\mathcal{L}\left(
A_{\ast}\left(  \gamma\right)  ,\lambda_{\ast}\left(  \gamma\right)
,\gamma\right)  =E_{U}\left(  A_{\ast}\left(  \gamma\right)  \right)
-\left\langle \left.  \lambda\right\vert \left\{  \Xi_{2}\left(  A_{\ast
}\left(  \gamma\right)  \right)  -\gamma\right\}  \right\rangle =E_{U}\left(
A_{\ast}\left(  \gamma\right)  \right)  .
\end{equation}
It can be shown \cite{Hestenes} that the sensitivity of the optimal
$N$-electron energy to changes in the density $\gamma$ is given by
\begin{equation}
\frac{\delta E_{U}\left(  A_{\ast}\left(  \gamma\right)  \right)  }%
{\delta\gamma}=\lambda_{\ast}\left(  \gamma\right)  .
\end{equation}
By definition
\begin{equation}
\mathcal{F}_{HK}\left(  \gamma\right)  =E_{U}\left(  A_{\ast}\left(
\gamma\right)  \right)
\end{equation}
which gives
\begin{equation}
\frac{\delta\mathcal{F}_{HK}\left(  \gamma\right)  }{\delta\gamma}%
=\lambda_{\ast}\left(  \gamma\right)  .
\end{equation}
One could find the ground state of the universal system by finding
\begin{equation}
\mathcal{F}_{HK}\left(  \gamma_{0}\right)  =\min_{\gamma\in\mathcal{P}%
}\mathcal{F}_{HK}\left(  \gamma\right)  , \label{EUniGS}%
\end{equation}
physically this would be the ground state of a uniform gas of electrons. At
such a point%
\begin{align}
&  \widehat{\left\{  K+V_{2}-\Omega\left(  \lambda_{\ast}\left(  \gamma
_{0}\right)  \right)  \right\}  }\left\vert A_{\ast}\left(  \gamma_{0}\right)
\right)  =0\nonumber\\
&  \mathcal{F}_{HK}\left(  \gamma_{0}\right)  =E_{U}\left(  A_{\ast}\left(
\gamma_{0}\right)  \right)  =\left\langle \left.  \lambda_{\ast}\left(
\gamma_{0}\right)  \right\vert \gamma_{0}\right\rangle .
\end{align}
The constraining potential $\Omega\left(  \lambda_{\ast}\left(  \gamma
_{0}\right)  \right)  $ is probably a uniform potential produced by a
background of positive charges. One can make eq. $\left(  \ref{EUniGS}\right)
$ an unconstrained optimization by expressing the density as
\begin{equation}
\gamma\left(  \mathbf{x}\right)  =\frac{Nf\left(  \mathbf{x}\right)  ^{2}%
}{\left\langle f|f\right\rangle }=\frac{N\left\langle f\left\vert \Phi\left(
\mathbf{x}\right)  ^{\dagger}\Phi\left(  \mathbf{x}\right)  f\right.
\right\rangle }{\left\langle f|f\right\rangle },\,f\in L_{2}\left(
\mathbb{X}\right)  ,
\end{equation}
giving%
\begin{equation}
\frac{\delta\mathcal{F}_{HK}\left(  \gamma\right)  }{\delta f}=\frac
{\delta\mathcal{F}_{HK}\left(  \gamma\right)  }{\delta\gamma}\frac
{\delta\gamma}{\delta f}=\int\frac{\delta\mathcal{F}_{HK}\left(
\gamma\right)  }{\delta\gamma\left(  \mathbf{x}\right)  }\frac{\delta
\gamma\left(  \mathbf{x}\right)  }{\delta f\left(  \mathbf{x}^{\prime}\right)
}d\mathbf{x}=\int\lambda_{\ast}\left(  \gamma\right)  \left(  \mathbf{x}%
\right)  \frac{\delta\gamma\left(  \mathbf{x}\right)  }{\delta f\left(
\mathbf{x}^{\prime}\right)  }d\mathbf{x,}%
\end{equation}
which can be evaluated by noting that%
\begin{equation}
\frac{\delta\gamma\left(  \mathbf{x}\right)  }{\delta f\left(  \mathbf{x}%
^{\prime}\right)  }=\frac{2N}{\left\langle f|f\right\rangle }\left\{
\left\langle f\left\vert \overset{}{\left\{  \Phi\left(  \mathbf{x}\right)
^{\dagger}\Phi\left(  \mathbf{x}\right)  -\gamma\left(  \mathbf{x}\right)
I\right\}  }\right.  \mathbf{x}^{\prime}\right\rangle \right\}  =\frac
{2N}{\left\langle f|f\right\rangle }\left\{  f\left(  \mathbf{x}\right)
\delta\left(  \mathbf{x-x}^{\prime}\right)  -\gamma\left(  \mathbf{x}\right)
f\left(  \mathbf{x}^{\prime}\right)  \right\}
\end{equation}
and hence
\begin{align}
\frac{\delta\mathcal{F}_{HK}\left(  \gamma\right)  }{\delta f\left(
\mathbf{x}^{\prime}\right)  }  &  =\frac{2N}{\left\langle f|f\right\rangle
}\int\lambda_{\ast}\left(  \gamma\right)  \left(  \mathbf{x}\right)  \left\{
\left\langle f\left\vert \overset{}{\Phi\left(  \mathbf{x}\right)  ^{\dagger
}\Phi\left(  \mathbf{x}\right)  }\right.  \mathbf{x}^{\prime}\right\rangle
-\gamma\left(  \mathbf{x}\right)  f\left(  \mathbf{x}^{\prime}\right)
\right\}  d\mathbf{x}\nonumber\\
&  =\frac{2N}{\left\langle f|f\right\rangle }f\left(  \mathbf{x}^{\prime
}\right)  \left\{  \lambda_{\ast}\left(  \gamma\right)  \left(  \mathbf{x}%
^{\prime}\right)  -\int\lambda_{\ast}\left(  \gamma\right)  \left(
\mathbf{x}\right)  \gamma\left(  \mathbf{x}\right)  d\mathbf{x}\right\}
\nonumber\\
&  =\frac{2N}{\left\langle f|f\right\rangle }f\left(  \mathbf{x}^{\prime
}\right)  \left\{  \overset{}{\lambda_{\ast}\left(  \gamma\right)  \left(
\mathbf{x}^{\prime}\right)  -\left\langle \left.  \lambda_{\ast}\left(
\gamma\right)  \right\vert \gamma\right\rangle }\right\}  .
\end{align}
At a stationary point
\begin{equation}
\frac{\delta\mathcal{F}_{HK}\left(  \gamma\right)  }{\delta f}%
=0\Leftrightarrow\lambda_{\ast}\left(  \gamma\right)  \left(  \mathbf{x}%
^{\prime}\right)  -\left\langle \left.  \lambda_{\ast}\left(  \gamma\right)
\right\vert \gamma\right\rangle =0\,\forall\mathbf{x}^{\prime}, \label{Ustat}%
\end{equation}
i.e. $\lambda_{\ast}\left(  \gamma\right)  $ is constant. A special case of
this is $\lambda_{\ast}\left(  \gamma\right)  =0$ $=\frac{\delta
\mathcal{F}_{HK}\left(  \gamma\right)  }{\delta\gamma}.$The equality
$\lambda_{\ast}\left(  \gamma\right)  =0$ shows that the density constraint is
not active at such stationary points, while $\lambda_{\ast}\left(
\gamma\right)  =$ const. could correspond to a uniform positive background.

\subsubsection{Decomposition of the Hohenburg-Kohn Functional}

If a minimal point $D_{\ast}\left(  \gamma\right)  $ exists and is unique and
then one can define kinetic energy, coulomb and exchange-correlation
functionals by eq. $\left(  \ref{HK}\right)  $ as $.$
\begin{align}
\mathcal{F}_{K}\left(  \gamma\right)   &  =\operatorname{Tr}\left\{  KD_{\ast
}\left(  \gamma\right)  \right\}  ,\nonumber\\
\mathcal{F}_{C}\left(  \gamma\right)   &  =\int\frac{\gamma\left(
\mathbf{r,\xi}\right)  \gamma\left(  \mathbf{r}^{\prime}\mathbf{,\xi}\right)
}{\left\Vert \mathbf{r-r}^{\prime}\right\Vert }d\mathbf{r}d\mathbf{r}^{\prime
}d\mathbf{\xi},\nonumber\\
\mathcal{F}_{XC}\left(  \gamma\right)   &  =\operatorname{Tr}\left\{
V_{2}D_{\ast}\left(  \gamma\right)  \right\}  -\mathcal{F}_{C}\left(
\gamma\right)  ,
\end{align}
producing the decomposition
\begin{equation}
\mathcal{F}_{HK}\left(  \gamma\right)  =\mathcal{F}_{K}\left(  \gamma\right)
+\mathcal{F}_{C}\left(  \gamma\right)  +\mathcal{F}_{XC}\left(  \gamma\right)
.
\end{equation}
As the concept of an exchange-correlation functional is central to the theory
the existence of a unique $D_{\ast}\left(  \gamma\right)  $ is crucial.
Different degenerate $D_{\ast}\left(  \gamma\right)  ^{\prime}$s could give
different $\mathcal{F}_{XC}\left(  \gamma\right)  ^{\prime}$s and
$\mathcal{F}_{K}\left(  \gamma\right)  ^{\prime}$s.

\subsubsection{Molecular System Energy}

The total energy functional of a system that contains $N$ electrons and
experiences an environment described by a \emph{local} potential $\mathrm{v}$
is
\begin{equation}
E\left(  \gamma\mathbf{,}\mathrm{v}\right)  =\mathcal{F}_{HK}\left(
\gamma\right)  +f_{\gamma}\left(  \mathrm{v}\right)  ,
\end{equation}
where the explicit effect of the environment is described by the functional
\begin{equation}
f_{\gamma}\left(  \mathrm{v}\right)  =\int\mathrm{v}\left(  \mathbf{r,\xi
}\right)  \gamma\left(  \mathbf{r,\xi}\right)  d\mathbf{r}d\mathbf{\xi.}%
\end{equation}
The ground state energy $E_{GS}\left(  \mathrm{v}\right)  $ and the
corresponding space-spin density $\gamma_{GS}$ $\left(  \mathrm{v}\right)  $
are then determined by
\begin{equation}
E_{GS}\left(  \mathrm{v}\right)  =E\left(  \gamma_{GS}\mathbf{,}%
\mathrm{v}\right)  =\min_{\mathbf{\gamma\in}\mathcal{D}}E\left(
\gamma\mathbf{,}\mathrm{v}\right)  ,
\end{equation}
where
\begin{equation}
\mathcal{D}=\left\{  \gamma\left\vert \int\gamma\left(  \mathbf{r,\xi}\right)
d\mathbf{r}d\mathbf{\xi}=N,\quad\gamma\left(  \mathbf{r,\xi}\right)
\geq0,\,\gamma\in L_{1}\left(  \mathbb{X}\right)  \right.  \right\}  .
\end{equation}
Clearly
\begin{equation}
\widehat{\left\{  K+V_{2}+\mathrm{v}-\Omega\left(  \lambda_{\ast}\left(
\gamma\right)  \right)  \right\}  }\left\vert A_{\ast}\left(  \gamma\right)
\right)  =0 \label{vstat}%
\end{equation}
as $f_{\gamma}\left(  \mathrm{v}\right)  $ is just a constant for fixed
$\gamma$, which can be recast as
\begin{equation}
\widehat{\left\{  K+V_{2}-\Omega\left(  \lambda_{\ast}\left(  \gamma\right)
-\mathrm{v}\right)  \right\}  }\left\vert A_{\ast}\left(  \gamma\right)
\right)  =0,
\end{equation}
where
\begin{equation}
\Omega\left(  \lambda_{\ast}\left(  \gamma\right)  -\mathrm{v}\right)
=\int\left\{  \lambda_{\ast}\left(  \gamma\right)  -\mathrm{v}\right\}
\left(  \mathbf{x}\right)  \Phi\left(  \mathbf{x}\right)  ^{\dagger}%
\Phi\left(  \mathbf{x}\right)  d\mathbf{x.}%
\end{equation}
This shows that if there exists a $\left\vert A_{\ast}\left(  \gamma\right)
\right)  $ such that eq. $\left(  \ref{Stateig}\right)  $ is satisfied and
$\lambda_{\ast}\left(  \gamma\right)  -\mathrm{v}$ is invertible on
$\operatorname{Supp}\gamma$ then $\gamma$ is V-representable \cite{weiner2005}.

In order to find the ground state energy $E\left(  \ \gamma_{GS}%
\mathbf{,}\mathrm{v}\right)  $ one needs to solve a constrained optimization
problem, however by defining
\begin{equation}
\gamma\left(  \mathbf{x}\right)  =\frac{Nf\left(  \mathbf{x}\right)  ^{2}%
}{\left\langle f|f\right\rangle }=\frac{N\left\langle f\left\vert \Phi\left(
\mathbf{x}\right)  ^{\dagger}\Phi\left(  \mathbf{x}\right)  f\right.
\right\rangle }{\left\langle f|f\right\rangle },\,f\in L_{2}\left(
\mathbb{X}\right)  ,
\end{equation}
this becomes the unconstrained problem
\begin{equation}
\min_{f}E\left(  f\mathbf{,}\mathrm{v}\right)  ;\quad f\in L_{2}\left(
\mathbb{X}\right)  .
\end{equation}
The stationary points of $E\left(  f\mathbf{,}\mathrm{v}\right)  $ with
respect to variations of the density occur when
\begin{equation}
\frac{\delta E\left(  f\mathbf{,}\mathrm{v}\right)  }{\delta f}=\frac{\delta
E\left(  f\mathbf{,}\mathrm{v}\right)  }{\delta\gamma}\frac{\delta\gamma
}{\delta f}=\left\{  \frac{\delta\mathcal{F}_{HK}\left(  \gamma\right)
}{\delta\gamma}+\mathrm{v}\right\}  \frac{\delta\gamma}{\delta f}=0.
\end{equation}
As
\begin{equation}
\frac{\delta E\left(  f\mathbf{,}\mathrm{v}\right)  }{\delta f\left(
\mathbf{x}^{\prime}\right)  }=\frac{2N}{\left\langle f|f\right\rangle
}f\left(  \mathbf{x}^{\prime}\right)  \left\{  \overset{}{\lambda_{\ast
}\left(  \gamma\right)  \left(  \mathbf{x}^{\prime}\right)  -\left\langle
\left.  \lambda_{\ast}\left(  \gamma\right)  \right\vert \gamma\right\rangle
}\right\}
\end{equation}
one observes
\begin{equation}
\frac{\delta E\left(  f\mathbf{,}\mathrm{v}\right)  }{\delta f}%
=0\Leftrightarrow\left[  \lambda_{\ast}\left(  \gamma\right)  -\mathrm{v}%
\right]  \left(  \mathbf{x}^{\prime}\right)  -\left\langle \left.  \left[
\lambda_{\ast}\left(  \gamma\right)  -\mathrm{v}\right]  \right\vert
\gamma\right\rangle =0\,\forall\mathbf{x}^{\prime},
\end{equation}
showing that $\lambda_{\ast}\left(  \gamma\right)  -\mathrm{v}=$ const. at
stationary points. In particular at minimizers $\gamma_{GS}$ of $E\left(
\gamma\mathbf{,}\mathrm{v}\right)  $
\begin{equation}
\mathrm{v}=-\frac{\delta\mathcal{F}_{HK}\left(  \ \gamma_{GS}\right)  }%
{\delta\gamma}+\text{const. and }\Omega\left(  \lambda_{\ast}\left(
\gamma_{GS}\right)  -\mathrm{v}+\text{const.}\right)  =0,
\end{equation}
where const. $=\left\langle \left.  \lambda_{\ast}\left(  \gamma\right)
\right\vert \gamma\right\rangle =E_{GS}\left(  \mathrm{v}\right)  .$ This last
identity shows that the optimization problem is not constrained by $\gamma$,
which is consistent with finding the ground state of $\widehat{K+V_{2}%
+\mathrm{v}-E_{GS}\left(  \mathrm{v}\right)  I}$ over all states $\left\vert
A\right)  .$

\subsection{Kohn-Sham DFT}

\subsubsection{Correspondence between Orbitals and Densities}

The density $\gamma_{GS}$ $\left(  \mathrm{v}\right)  $ can be always be
expressed in terms of the Natural General Spin Orbitals (NGSO's) $\left\{
\psi_{j}\right\}  $ and occupation numbers $\left\{  N_{j}\right\}  $ of the
FORDO $D_{\ast}^{1}\left(  \gamma\right)  $ as%

\begin{equation}
\left[  \gamma_{GS}\left(  \mathrm{v}\right)  \right]  \left(  \mathbf{x}%
\right)  =\left[  D_{\ast}^{1}\left(  \gamma_{GS}\left(  \mathrm{v}\right)
\right)  \right]  \left(  \mathbf{x,x}\right)  =\sum_{1\leq j\leq r_{1}}%
N_{j}\left\vert \psi_{j}\left(  \mathbf{x}\right)  \right\vert ^{2},
\label{NGSOexpan}%
\end{equation}
but many other sets of GSO's also produce the same density, in particular
Harimanns \cite{Harimann} theorem ensures that there are always expansions of
the form
\begin{equation}
\left[  \gamma_{GS}\left(  \mathrm{v}\right)  \right]  \left(  \mathbf{x}%
\right)  =\left[  D_{\ast}^{1}\left(  \gamma_{GS}\left(  \mathrm{v}\right)
\right)  \right]  \left(  \mathbf{x,x}\right)  =\sum_{1\leq j\leq N}\left\vert
\varphi_{j}\left(  \mathbf{x}\right)  \right\vert ^{2} \label{Genexp}%
\end{equation}
in terms of square summable functions $\left\{  \left.  \varphi_{j}\right\vert
1\leq j\leq N\right\}  $ for any density $\gamma$.

In general the density $\gamma$ is a functional of orthogonal sets of GSO's
$\mathbf{\psi}$
\begin{equation}
\left\langle \left.  \psi_{j}\right\vert \psi_{k}\right\rangle =N_{j}%
\delta_{jk},0\leq N_{j}\leq1,1\leq j,k\leq r_{1},
\end{equation}
defined by%
\begin{align}
\gamma &  =\Gamma\left(  \mathbf{\psi}\right)  ,\nonumber\\
\left[  \Gamma\left(  \mathbf{\psi}\right)  \right]  \left(  \mathbf{x}%
\right)   &  =\sum_{1\leq j\leq r_{1}}\left\vert \psi_{j}\left(
\mathbf{x}\right)  \right\vert ^{2}=\gamma\left(  \mathbf{x}\right)  ,
\end{align}
where
\begin{equation}
\Gamma:B_{1}\left(  \mathcal{H}^{1}\right)  \subset\mathcal{H}^{1}\rightarrow
L_{1+}\left(  \mathbb{X}\right)  \subset L_{1}\left(  \mathbb{X}\right)
\end{equation}
and $B_{1}\left(  \mathcal{H}^{1}\right)  $ is the unit ball of $\mathcal{H}%
^{1}$ i.e.
\begin{equation}
B_{1}\left(  \mathcal{H}^{1}\right)  =\left\{  \left.  \psi\right\vert
\left\Vert \psi\right\Vert \leq1\right\}  .
\end{equation}
In the preceeding definition some of the functions may be identically zero. As
many sets of GSO's $\mathbf{\psi}$ produce the same density the inverse
$\Gamma^{-1}$ does not exist. One could restrict attention to the real Hilbert
space $L_{2}\left(  \mathbb{X}\right)  $ instead of $\mathcal{H}^{1}$ and use
real valued functions in the expansions eqs.$\left(  \ref{NGSOexpan}\right)  $
and $\left(  \ref{Genexp}\right)  $ e.g.
\begin{equation}
\left[  \gamma_{GS}\left(  \mathrm{v}\right)  \right]  \left(  \mathbf{x}%
\right)  =\left[  D_{\ast}^{1}\left(  \gamma_{GS}\left(  \mathrm{v}\right)
\right)  \right]  \left(  \mathbf{x,x}\right)  =\sum_{1\leq j\leq r_{1}}%
N_{j}\tilde{\psi}_{j}\left(  \mathbf{x}\right)  ^{2}=\sum_{1\leq j\leq r_{1}%
}\varphi_{j}\left(  \mathbf{x}\right)  ^{2},
\end{equation}
where %

\begin{align}
\psi_{j}\left(  \mathbf{x}\right)    & =\tilde{\psi}_{j}\left(  \mathbf{x}%
\right)  e^{i\theta_{j}\left(  \mathbf{x}\right)  }\nonumber\\
\varphi_{j}\left(  \mathbf{x}\right)    & =N_{j}^{\frac{1}{2}}\psi_{j}\left(
\mathbf{x}\right)  ,
\end{align}
but one looses some physical information e.g. the expectation value of the
kinetic operator. Of course the expansion eq. $\left(  \ref{NGSOexpan}\right)
$ provides a unique correspondence up to unitary mixing of NGSO's with same
occupation number, but one has to know the FORDO of the exact $N$-electron
ground state in order use it.

A more accessible route to the production an inverse function $\Gamma^{-1}$ is
to consider a universal kinetic energy functional $K_{U}$ of the $N$-electron
state and a kinetic energy functional $\mathcal{F}_{T}$ of the density defined
by
\begin{equation}
\mathcal{F}_{T}\left(  \gamma\right)  =K_{U}\left(  D_{\#}\left(
\gamma\right)  \right)  =\operatorname{Tr}\left\{  KD_{\#}\left(
\gamma\right)  \right\}  =\operatorname{Tr}\left\{  KA_{\#}^{2}\left(
\gamma\right)  \right\}  =\min_{D\in\left[  \gamma\right]  _{N}}%
\operatorname{Tr}\left\{  KA_{\#}^{2}\left(  \gamma\right)  \right\}  .
\label{KEFunc}%
\end{equation}
If a bound state minimum exists for some local one body potential one has
\begin{equation}
\widehat{\left\{  K-\Omega\left(  \lambda_{\#}\left(  \gamma\right)  \right)
\right\}  }\left\vert A_{\#}\left(  \gamma\right)  \right)  =0,
\end{equation}
where $A_{\#}\left(  \gamma\right)  \in\mathcal{B}_{2sa}\left(  \mathcal{H}%
^{N}\right)  $ and $\left(  A_{\#}\left(  \gamma\right)  ,\lambda_{\#}\left(
\gamma\right)  \right)  $ is an unconstrained saddle point of the Lagrangian
\begin{equation}
\mathcal{L}_{K}\left(  A,\lambda,\gamma\right)  =K_{U}\left(  A\right)
-\Omega_{A}\left(  \lambda\right)  +\left\langle \left.  \lambda\right\vert
\gamma\right\rangle =\left(  A\left\vert \left\{  \widehat{\left\{
K-\Omega\left(  \lambda\right)  \right\}  }+\left\langle \left.
\lambda\right\vert \gamma\right\rangle \hat{I}\right\}  \right.  A\right)  .
\end{equation}
This gives a stationary state $D_{\#}\left(  \gamma\right)  =A_{\#}\left(
\gamma\right)  ^{2}$ of $\left\{  K-\Omega\left(  \lambda_{\#}\left(
\gamma\right)  \right)  \right\}  $ i.e.
\begin{equation}
\left[  \left\{  K-\Omega\left(  \lambda_{\#}\left(  \gamma\right)  \right)
\right\}  ,D_{\#}\left(  \gamma\right)  \right]  =0,
\end{equation}
which in turn defines a FORDO $D_{\#}^{1}\left(  \gamma\right)  $ given by
(see Appendix \ref{Denmat})
\begin{equation}
D_{\#}^{1}\left(  \gamma\right)  =NL_{N}^{1}\left(  D_{\#}\left(
\gamma\right)  \right)  .
\end{equation}
As $K-\Omega\left(  \lambda_{\#}\left(  \gamma\right)  \right)  $ is an one
electron operator the state $D_{\#}\left(  \gamma\right)  $ must be a mixture
of IPS's or a pure IPS and thus determined by the eigenvalues $\left\{
N_{\#j}\left(  \gamma\right)  \right\}  $ and eigenvectors $\left\{
\left\vert \widetilde{\varphi_{\#j}\left(  \gamma\right)  }\right\rangle
\right\}  $of $D_{\#}^{1}\left(  \gamma\right)  $. Expanding $D_{\#}%
^{1}\left(  \gamma\right)  $ in terms of normalized NGSO's $\left\{  \left.
\left\vert \widetilde{\varphi_{\#j}}\right\rangle \right\vert 1\leq j\leq
r\right\}  $ gives
\begin{equation}
D_{\#}^{1}\left(  \gamma\right)  =\sum_{1\leq j\leq r_{1}}N_{\#j}\left(
\gamma\right)  \left\vert \widetilde{\varphi_{\#j}\left(  \gamma\right)
}\right\rangle \left\langle \widetilde{\varphi_{\#j}\left(  \gamma\right)
}\right\vert ,
\end{equation}
which after renormalizing the NGSO's to the occupation numbers i.e.
\begin{equation}
\left\vert \varphi_{\#j}\left(  \gamma\right)  \right\rangle =N_{\#j}\left(
\gamma\right)  ^{\frac{1}{2}}\left\vert \widetilde{\varphi_{\#j}\left(
\gamma\right)  }\right\rangle ,1\leq j\leq r_{1}%
\end{equation}
results in
\begin{equation}
\gamma\left(  \mathbf{x}\right)  =\sum_{1\leq j\leq r_{1}}\left\vert \left[
\varphi_{\#j}\left(  \gamma\right)  \right]  \left(  \mathbf{x}\right)
\right\vert ^{2}. \label{gammaKS}%
\end{equation}
If the constrained minimum of $K_{U}$ is non degenerate then there is a unique
correspondence, up to unitary mixing of degenerate eigenvectors of $D_{\#}%
^{1}\left(  \gamma\right)  ,$ between $\gamma$ and $\left\{  \left\vert
\mathbf{\varphi}_{\#}\right\rangle \right\}  ,$ thus modulo this unitary
mixing
\begin{equation}
\mathbf{\varphi}_{\#}=\Gamma_{\#}^{-1}\left(  \gamma\right)  .
\end{equation}
If the constrained minimum of $K_{U}$ is degenerate, e.g. corresponding to
FORDO's $\left\{  \left.  D_{\#\kappa}^{1}\left(  \gamma\right)  \right\vert
1\leq\kappa\leq m\right\}  $ of the stationary states $\left\{  \left.
D_{\#\kappa}\left(  \gamma\right)  \right\vert 1\leq\kappa\leq m\right\}  $ of
$K-\Omega\left(  \lambda_{\#}\left(  \gamma\right)  \right)  $ then as
\begin{equation}
\left[  \left\{  K-\Omega\left(  \lambda_{\#}\left(  \gamma\right)  \right)
\right\}  ,D_{\#\kappa}\left(  \gamma\right)  \right]  =\left[  \left\{
K-\Omega\left(  \lambda_{\#}\left(  \gamma\right)  \right)  \right\}
,D_{\#\kappa}^{1}\left(  \gamma\right)  \right]  =0,
\end{equation}
the space of degenerate states corresponds to the same set of NGSO's $\left\{
\left\vert \mathbf{\varphi}_{\#}\left(  \gamma\right)  \right\rangle \right\}
$ but normalized to different occupation numbers $\left\{  \mathbf{N}%
_{\#}\left(  \gamma\right)  \right\}  ,$ where $\mathbf{N}_{\#}\left(
\gamma\right)  $ are convex combinations of $\left\{  \left.  \mathbf{N}%
_{\#\kappa}\left(  \gamma\right)  \right\vert 1\leq\kappa\leq m\right\}  $.
Hence if degeneracies occur these are reflected in multiple optimal values for
$\mathbf{\mathbf{N}_{\#}}$ for fixed $\left\{  \left\vert \widetilde
{\mathbf{\varphi}_{\#}\left(  \gamma\right)  }\right\rangle \right\}
\mathbf{.}$

The values of the universal kinetic energy functional $\mathcal{F}_{T}\left(
\gamma\right)  $ also correspond to the minima of the following constrained
problem, expressed in terms of $\left\{  \left.  N_{j},\left\vert \varphi
_{j}\right\rangle \right\vert 1\leq j\leq r_{1}\right\}  $ as the sequential
minimization
\begin{subequations}
\begin{equation}
\mathcal{F}_{T}\left(  \mathbf{N}_{\#}\mathbf{,\varphi}_{\#}\right)
=\mathcal{F}_{T}\left(  \gamma\left(  \mathbf{N}_{\#}\mathbf{,\varphi}%
_{\#}\right)  \right)  =K_{U}\left(  \mathbf{N}_{\#}\mathbf{,\varphi}%
_{\#},\gamma\right)  =\min_{\mathbf{N,\varphi}}K_{U}\left(  \mathbf{N,\varphi
},\gamma\right)  =\min_{\mathbf{N}}\min_{\mathbf{\varphi}}\left\{  \sum_{1\leq
j\leq r_{1}}\left\langle \varphi_{j}\left\vert K\varphi_{j}\right.
\right\rangle \right\}  ,
\end{equation}
where $\gamma$ and $\mathbf{N}$ play the role of control parameters and the
first minimization is constrained by
\end{subequations}
\begin{align}
\quad &  \sum_{1\leq j\leq r_{1}}\left\langle \varphi_{j}\left\vert
\Phi\left(  \mathbf{x}\right)  ^{\dagger}\Phi\left(  \mathbf{x}\right)
\right.  \varphi_{j}\right\rangle =\gamma\left(  \mathbf{x}\right)
,\nonumber\\
&  \int\gamma\left(  \mathbf{x}\right)  d\mathbf{x}=N,\nonumber\\
&  \left\langle \left.  \varphi_{j}\right\vert \varphi_{k}\right\rangle
=N_{j}\delta_{jk},\,1\leq j,k\leq r_{1}%
\end{align}
and the second by
\begin{equation}
0\leq N_{j}\leq1,\,1\leq j\leq r_{1}.
\end{equation}
The Lagrangian for the first minimization, where $\mathbf{N}$ and $\gamma$ are
fixed is%
\[
\mathcal{L}\left(  \mathbf{\varphi},\lambda,\mathbf{\nu},\gamma,\mathbf{N}%
\right)  =K_{U}\left(  \mathbf{\varphi}\right)  -\left\langle \left.
\lambda\right\vert \left\{  \Xi_{IP}\left(  \mathbf{\varphi}\right)
-\gamma\right\}  \right\rangle -\sum_{1\leq j,k\leq r_{1}}\nu_{jk}\left\{
N_{j}\delta_{jk}-\left\langle \left.  \varphi_{j}\right\vert \varphi
_{k}\right\rangle \right\}  ,
\]
where
\begin{equation}
\Xi_{IP}\left(  \mathbf{\varphi}\right)  =\sum_{1\leq j\leq r}\left\langle
\varphi_{j}\left\vert \Phi\left(  \mathbf{x}\right)  ^{\dagger}\Phi\left(
\mathbf{x}\right)  \right.  \varphi_{j}\right\rangle
\end{equation}
and hence
\begin{equation}
\sum_{1\leq j\leq r_{1}}\left\langle \varphi_{j}\left\vert \Omega\left(
\lambda\right)  \right.  \varphi_{j}\right\rangle =\left\langle \left.
\lambda\right\vert \left\{  \Xi_{IP}\left(  \mathbf{\varphi}\right)
-\gamma\right\}  \right\rangle .
\end{equation}
The solution of this optimization is
\begin{align}
K_{U}\left(  \mathbf{\varphi}_{\#}\left(  \gamma,\mathbf{N}\right)  \right)
&  =\mathcal{L}\left(  \mathbf{\varphi}_{\#},\lambda_{\#},\mathbf{\nu}%
_{\#},\gamma,\mathbf{N}\right)  =\max_{\lambda,\mathbf{\nu}}\min
_{\mathbf{\varphi}}\mathcal{L}\left(  \mathbf{\varphi},\lambda,\mathbf{\nu
},\gamma,\mathbf{N}\right) \nonumber\\
&  =\max_{\lambda,\mathbf{\nu}}\min_{\mathbf{\varphi}}\left\{  \sum_{1\leq
j\leq r_{1}}\left\langle \varphi_{j}\left\vert \left\{  \overset{}%
{K-\Omega\left(  \lambda\right)  }\right\}  \varphi_{j}\right.  \right\rangle
-\sum_{1\leq j,k\leq r_{1}}\nu_{jk}\left\{  N_{j}\delta_{jk}-\left\langle
\left.  \varphi_{j}\right\vert \varphi_{k}\right\rangle \right\}  \right\}  .
\end{align}
Sensitivity theory gives that at optimal $\mathbf{\varphi}_{\#}$'s,
$\nu_{\#jj}$'s and $\lambda_{\#}$'s
\begin{equation}
\frac{\partial K_{U}\left(  \mathbf{\varphi}_{\#}\left(  \gamma,\mathbf{N}%
\right)  \right)  }{\partial N_{j}}=\nu_{\#jj}\left(  \mathbf{N}\right)
\end{equation}
and
\begin{equation}
\frac{\delta K_{U}\left(  \mathbf{\varphi}_{\#}\left(  \gamma,\mathbf{N}%
\right)  \right)  }{\delta\gamma\left(  \mathbf{x}\right)  }=\left[
\lambda_{\#}\left(  \gamma\right)  \right]  \left(  \mathbf{x}\right)  .
\end{equation}
One can then vary the control parameter $\mathbf{N}$ in an unconstrained
fashion, while keeping $\gamma$ fixed (thus keeping $\sum_{1\leq j,k\leq
r_{1}}N_{j}=N,$ by setting
\begin{equation}
N_{j}=\sin^{2}\theta_{j},1\leq\theta_{j}\leq2\pi,1\leq j\leq r_{1}.
\end{equation}
This gives
\begin{equation}
\frac{\partial\mathcal{L}\left(  \mathbf{\varphi},\lambda,\mathbf{\nu}%
,\gamma,\mathbf{N}\right)  }{\partial\theta_{l}}=\sum_{1\leq j\leq r_{1}}%
\frac{\partial\mathcal{L}\left(  \mathbf{\varphi},\lambda,\mathbf{\nu}%
,\gamma,\mathbf{N}\right)  }{\partial N_{j}}\sin2\theta_{l}\delta_{jl}%
=\frac{\partial\mathcal{L}\left(  \mathbf{\varphi},\lambda,\mathbf{\nu}%
,\gamma,\mathbf{N}\right)  }{\partial N_{l}}\sin2\theta_{l}.
\end{equation}
The universal kinetic energy functional can then be defined as
\begin{equation}
\mathcal{F}_{T}\left(  \mathbf{\varphi}_{\#}\left(  \gamma,\mathbf{N}%
_{\#}\right)  \right)  =\mathcal{F}_{T}\left(  \Gamma^{-1}\left(
\gamma\right)  \right)  =\mathcal{F}_{T}\left(  \gamma\right)  =K_{U}\left(
\mathbf{\varphi}_{\#}\left(  \gamma,\mathbf{N}_{\#}\right)  \right)
=\min_{\mathbf{N}}\min_{\mathbf{\varphi}}K_{U}\left(  \mathbf{\varphi}\left(
\gamma,\mathbf{N}\right)  \right)
\end{equation}
and can be expanded as
\begin{equation}
\mathcal{F}_{T}\left(  \gamma\right)  =\sum_{1\leq j\leq r_{1}}\left\langle
\widetilde{\varphi_{\#j}\left(  \gamma\right)  }\left\vert K\widetilde
{\varphi_{\#j}\left(  \gamma\right)  }\right.  \right\rangle N_{\#j}\left(
\gamma\right)  =\sum_{1\leq j\leq r_{1}}\left\langle \varphi_{\#j}\left(
\gamma,\mathbf{N}_{\#}\right)  \left\vert K\varphi_{\#j}\left(  \gamma
,\mathbf{N}_{\#}\right)  \right.  \right\rangle ,
\end{equation}
where there may be multiple values of $\mathbf{N}_{\#}$. By sensitivity theory
one has that for any value of the control parameter $\mathbf{N}$
\begin{equation}
\frac{\partial K_{U}\left(  \mathbf{\varphi}_{\#}\left(  \gamma,\mathbf{N}%
\right)  \right)  }{\partial N_{j}}=\nu_{\#jj}\left(  \mathbf{N}\right)  ,
\end{equation}
thus
\begin{equation}
\frac{\partial\mathcal{F}_{T}\left(  \mathbf{\varphi}_{\#}\left(
\gamma,\mathbf{N}\right)  \right)  }{\partial N_{j}}=\frac{\partial
K_{U}\left(  \mathbf{\varphi}_{\#}\left(  \gamma,\mathbf{N}\right)  \right)
}{\partial N_{j}}=\nu_{\#jj}\left(  \mathbf{N}\right)  .
\end{equation}
Thus at stationary points
\begin{equation}
\sin2\theta_{j}=0\text{ or }\frac{\partial\mathcal{L}\left(  \mathbf{\varphi
}_{\#},\lambda_{\#},\mathbf{\nu}_{\#},\gamma,\mathbf{N}_{\#}\right)
}{\partial N_{j}}=\frac{\partial K_{U}\left(  \mathbf{\varphi}_{\#}\left(
\gamma,\mathbf{N}\right)  \right)  }{\partial N_{j}}=\nu_{\#jj}\left(
\mathbf{N}\right)  =0,
\end{equation}
the first condition implies that
\begin{equation}
\theta_{j}=n\frac{\pi}{2}\Rightarrow\sin\theta_{j}=1\Rightarrow N_{j}=1.
\end{equation}
Thus at stationary points $\mathbf{N}_{\#}$, either $\nu_{\#jj}\left(
\mathbf{N}_{\#}\right)  =0$ or $N_{\#j}=1.$

\subsubsection{Model KS-IPS}

One can define for a given space-spin density $\gamma$ in an environment
characterized by a local potential $\mathrm{w}$ a total energy functional
\begin{equation}
\mathcal{E}_{NI}\left(  \gamma\mathbf{,}\mathrm{w}\right)  =\mathcal{F}%
_{T}\left(  \gamma\right)  +f_{\gamma}\left(  \mathrm{w}\right)  ,
\end{equation}
and then determine a local potential $\mathrm{w}_{\#}$ such that
\begin{equation}
E_{NIGS}\left(  \mathrm{w}_{\#}\right)  =\mathcal{E}_{NIGS}\left(  \gamma
_{GS}\left(  \mathrm{v}\right)  \mathbf{,}\mathrm{w}_{\#}\right)
=\min_{\mathbf{\gamma\in}\mathcal{D}}\mathcal{E}_{NI}\left(  \gamma
\mathbf{,}\mathrm{w}_{\#}\right)  . \label{NIGS}%
\end{equation}
As the Hamiltonian describing noninteracting electrons is a one body
Hamiltonian the ground state density operator $D_{\#}\left(  \gamma
_{GS}\left(  \mathrm{v}\right)  \mathbf{,}\mathrm{w}_{\#}\right)  $ must be an
ensemble or a pure Independent Particle State (IPS)\ density operator. Thus
the space-spin density of the ground state of the interacting system in an
local external potential $\mathrm{v}$ is the same as space-spin density of the
ground state of a non interacting system determined by an environment
$\mathrm{w}_{\#}\mathrm{.}$

We can express the non interacting energy in terms of the orthogoanl GSO
variables $\mathbf{\varphi\in}B_{1}\left(  L_{2}\left(  \mathbb{X}\right)
\right)  $ as
\begin{align}
&  \mathcal{E}_{NI}\left(  \Gamma\left(  \mathbf{\varphi}\right)
\mathbf{,}\mathrm{w}\right)  =\mathcal{F}_{T}\left(  \Gamma\left(
\mathbf{\varphi}\right)  \right)  +f_{\Gamma\left(  \mathbf{\varphi}\right)
}\left(  \mathrm{w}\right) \nonumber\\
&  \equiv\mathcal{E}_{NI}\left(  \mathbf{\varphi,}\mathrm{w}\right)
=\mathcal{F}_{T}\left(  \mathbf{\varphi}\right)  +f_{\mathbf{\varphi}}\left(
\mathrm{w}\right)  ,
\end{align}
where we have set $\mathcal{F}\circ\Gamma\equiv\mathcal{F}.$ Explicitly the
functional $\mathcal{E}_{NI}\left(  \mathbf{\varphi,}\mathrm{w}\right)  $ is
\begin{equation}
\mathcal{E}_{NI}\left(  \mathbf{\varphi,}\mathrm{w}\right)  =\sum_{1\leq j\leq
r_{1}}\left\langle \varphi_{j}\left\vert \left\{  K+\mathrm{w}\right\}
\varphi_{j}\right.  \right\rangle =\sum_{1\leq j\leq r_{1}}\left\langle
\tilde{\varphi}_{j}\left\vert \left\{  K+\mathrm{w}\right\}  \tilde{\varphi
}_{j}\right.  \right\rangle N_{j}, \label{NI}%
\end{equation}
where
\begin{equation}
\left\vert \varphi_{j}\right\rangle =N_{j}^{\frac{1}{2}}\left\vert
\tilde{\varphi}_{j}\right\rangle ;1\leq j\leq r_{1}%
\end{equation}
and
\begin{equation}
\left\langle \tilde{\varphi}_{j}\left\vert \tilde{\varphi}_{k}\right.
\right\rangle =\delta_{jk};\quad\left\langle \varphi_{j}\left\vert \varphi
_{k}\right.  \right\rangle =N_{j}\delta_{jk};\quad1\leq j,k\leq r_{1}.
\end{equation}
The minimization eq. $\left(  \ref{NIGS}\right)  $
\begin{equation}
E_{NIGS}\left(  \mathrm{w}_{\#}\right)  =\mathcal{E}_{NIGS}\left(
\mathbf{\varphi}_{GS}\left(  \mathrm{v}\right)  \mathbf{,}\mathrm{w}%
_{\#}\right)  =\min_{\mathbf{\varphi\in}\mathfrak{S}}\mathcal{E}_{NI}\left(
\mathbf{\varphi,}\mathrm{w}_{\#}\right)  =\min_{\mathbf{\varphi\in
}\mathfrak{\tilde{S}}}\mathcal{E}_{NI}\left(  \mathbf{\tilde{\varphi}%
,}\mathrm{w}_{\#}\right)  \label{MNIGS}%
\end{equation}
can be expressed equivalently in either of the two forms of eq. $\left(
\ref{NI}\right)  $ i.e. in terms of $\mathbf{\varphi}$ normalized to
$\mathbf{N}$ or the orthonormal set $\mathbf{\tilde{\varphi}}$ with
$\mathbf{N}$ being explicitly expressed in the functional $\mathcal{E}_{NI}.$
then becomes where the set of feasable values is defined as
\begin{align}
\mathfrak{S}  &  =\left\{  \mathbf{\varphi}\left\vert \mathbf{\varphi\in}%
B_{1}\left(  L_{2}\left(  \mathbb{X}\right)  \right)  ;\,\int\left[
\Gamma\left(  \mathbf{\varphi}\right)  \right]  \left(  \mathbf{x}\right)
d\mathbf{x}=N;\,\left\langle \left.  \varphi_{j}\right\vert \varphi
_{k}\right\rangle =N_{j}\delta_{jk},1\leq j,k\leq r_{1}\right.  \right\}
\nonumber\\
&  \equiv\left\{  \mathbf{\varphi}\left\vert \left\langle \left.  \varphi
_{j}\right\vert \varphi_{k}\right\rangle =N_{j}\delta_{jk};\,0\leq N_{j}%
\leq1,1\leq j,k\leq r_{1};\,\sum_{1\leq j\leq r_{1}}N_{j}=N\right.  \right\}
\end{align}
and
\begin{align}
\mathfrak{\tilde{S}}  &  =\left\{  \mathbf{\tilde{\varphi}}\left\vert
\mathbf{\tilde{\varphi}\in}S_{1}\left(  L_{2}\left(  \mathbb{X}\right)
\right)  ;\,\int\left[  \tilde{\Gamma}\left(  \mathbf{\tilde{\varphi}}\right)
\right]  \left(  \mathbf{x}\right)  d\mathbf{x}=N;\,\left\langle \left.
\tilde{\varphi}_{j}\right\vert \tilde{\varphi}_{k}\right\rangle =\delta
_{jk},1\leq j,k\leq r_{1}\right.  \right\} \nonumber\\
&  \equiv\left\{  \mathbf{\tilde{\varphi}}\left\vert \left\langle \left.
\tilde{\varphi}_{j}\right\vert \tilde{\varphi}_{k}\right\rangle =\delta
_{jk};\,0\leq N_{j}\leq1,1\leq j,k\leq r_{1};\,\sum_{1\leq j\leq r_{1}}%
N_{j}=N\right.  \right\}  .
\end{align}
We will focus on the optimization in terms of $\mathbf{\varphi}$. The
minimization eq. $\left(  \ref{MNIGS}\right)  $ can be considered as a
sequential optimization
\begin{equation}
E_{NIGS}\left(  \mathrm{w}_{\#}\right)  =\min_{\mathbf{N\in}\mathcal{N}}%
\min_{\mathbf{\varphi\in}\mathfrak{S}_{\mathbf{N}}}\mathcal{E}_{NI}\left(
\mathbf{\varphi,}\mathrm{w}_{\#}\right)  , \label{seqopt}%
\end{equation}
where
\begin{equation}
\mathfrak{S}_{\mathbf{N}}=\left\{  \mathbf{\varphi}\left\vert \,\,\left\langle
\left.  \varphi_{j}\right\vert \varphi_{k}\right\rangle =N_{j}\delta
_{jk},1\leq j,k\leq r_{1}\right.  \right\}
\end{equation}
and
\begin{equation}
\mathcal{N}=\left\{  \mathbf{N}\left\vert 0\leq N_{j}\leq1,1\leq j\leq
r_{1};\sum_{1\leq j\leq r_{1}}N_{j}=N\right.  \right\}  .
\end{equation}
The Lagrangian for the first optimization is
\begin{equation}
\mathcal{L}_{K}\left(  \mathbf{\varphi},\mathbf{\varepsilon,N,}\mathrm{w}%
_{\#}\right)  =\mathcal{E}_{NI}\left(  \mathbf{\varphi,}\mathrm{w}%
_{\#}\right)  -\sum_{1\leq j,k\leq r_{1}}\varepsilon_{jk}\left(  \left\langle
\left.  \varphi_{j}\right\vert \varphi_{k}\right\rangle -N_{j}\delta
_{jk}\right)  ,
\end{equation}
which gives the constrained minimum $\mathcal{E}_{NI}\left(  \mathbf{\varphi
}\left(  \mathbf{N}\right)  \mathbf{,}\mathrm{w}_{\#}\right)  $ of
$\mathcal{E}_{NI}\left(  \mathbf{\varphi,}\mathrm{w}_{\#}\right)  $ as the
unconstrained saddle point
\begin{equation}
\mathcal{L}_{K}\left(  \mathbf{\varphi}\left(  \mathbf{N}\right)
,\mathbf{\varepsilon\left(  \mathbf{N}\right)  ,N,}\mathrm{w}_{\#}\right)
=\max_{\mathbf{\varepsilon}}\min_{\mathbf{\varphi\in}\mathfrak{S}_{\mathbf{N}%
}}\mathcal{L}_{K}\left(  \mathbf{\varphi},\mathbf{\varepsilon,N,}%
\mathrm{w}_{\#}\right)  .
\end{equation}
The necessary conditions for this saddle point are
\begin{subequations}
\begin{align}
\frac{\delta\mathcal{L}_{K}\left(  \mathbf{\varphi},\mathbf{\varepsilon
,N,}\mathrm{w}_{\#}\right)  }{\delta\varphi_{j}}  &  =0;\,1\leq j\leq
r_{1}\label{LK1}\\
\frac{\partial\mathcal{L}_{K}\left(  \mathbf{\varphi},\mathbf{\varepsilon
,N,}\mathrm{w}_{\#}\right)  }{\partial\varepsilon_{jk}}  &  =0;\,1\leq j,k\leq
r_{1}. \label{LK2}%
\end{align}
The condition eq. $\left(  \ref{LK2}\right)  $ just gives back the constraint,
while condition eq. $\left(  \ref{LK1}\right)  $ gives for the optimal point
$\mathbf{\varphi}\left(  \mathbf{N}\right)  $
\end{subequations}
\begin{equation}
\frac{\delta\mathcal{E}_{NI}\left(  \mathbf{\varphi}\left(  \mathbf{N}\right)
\mathbf{,}\mathrm{w}_{\#}\right)  }{\delta\varphi_{j}}-2\sum_{1\leq k\leq
r_{1}}\varepsilon_{jk}\left(  \mathbf{N}\right)  \left\vert \varphi_{k}\left(
\mathbf{N}\right)  \right\rangle =0;\,1\leq j\leq r_{1},
\end{equation}
which after noting that
\begin{align}
\frac{\delta\mathcal{E}_{NI}\left(  \mathbf{\varphi,}\mathrm{w}_{\#}\right)
}{\delta\varphi_{j}\left(  \mathbf{x}\right)  }  &  =\int\frac{\delta
\mathcal{E}_{NI}\left(  \gamma\mathbf{,}\mathrm{w}_{\#}\right)  }{\delta
\gamma\left(  \mathbf{x}^{\prime}\right)  }\frac{\delta\gamma\left(
\mathbf{x}^{\prime}\right)  }{\delta\varphi_{j}\left(  \mathbf{x}\right)
}d\mathbf{x}^{\prime}\nonumber\\
&  =2\int\frac{\delta\mathcal{E}_{NI}\left(  \gamma\mathbf{,}\mathrm{w}%
_{\#}\right)  }{\delta\gamma\left(  \mathbf{x}\right)  }\varphi_{j}\left(
\mathbf{x}^{\prime}\right)  \delta\left(  \mathbf{x-x}^{\prime}\right)
d\mathbf{x}^{\prime}=2\frac{\delta\mathcal{E}_{NI}\left(  \gamma
\mathbf{,}\mathrm{w}_{\#}\right)  }{\delta\gamma\left(  \mathbf{x}\right)
}\varphi_{j}\left(  \mathbf{x}\right)
\end{align}
gives%
\begin{equation}
\frac{\delta\mathcal{E}_{NI}\left(  \gamma\left(  \mathbf{\varphi}\left(
\mathbf{N}\right)  \right)  \mathbf{,}\mathrm{w}_{\#}\right)  }{\delta\gamma
}\left\vert \varphi_{j}\left(  \mathbf{N}\right)  \right\rangle -\sum_{1\leq
k\leq r_{1}}\varepsilon_{jk}\left(  \mathbf{N}\right)  \left\vert \varphi
_{k}\left(  \mathbf{N}\right)  \right\rangle =0;\,1\leq j\leq r_{1}.
\label{NIstat}%
\end{equation}
The operator $\frac{\delta\mathcal{E}_{NI}\left(  \gamma\mathbf{,}%
\mathrm{w}_{\#}\right)  }{\delta\gamma}$ is a multiplication operator given
by
\begin{equation}
\frac{\delta\mathcal{E}_{NI}\left(  \gamma\left(  \mathbf{\varphi}\left(
\mathbf{N}\right)  \right)  \mathbf{,}\mathrm{w}_{\#}\right)  }{\delta\gamma
}\left\vert \varphi_{k}\left(  \mathbf{N}\right)  \right\rangle =\left\vert
\frac{\delta\mathcal{E}_{NI}\left(  \gamma\left(  \mathbf{\varphi}\left(
\mathbf{N}\right)  \right)  \mathbf{,}\mathrm{w}_{\#}\right)  }{\delta\gamma
}\varphi_{k}\left(  \mathbf{N}\right)  \right\rangle ,
\end{equation}
where
\begin{equation}
\left[  \frac{\delta\mathcal{E}_{NI}\left(  \gamma\left(  \mathbf{\varphi
}\left(  \mathbf{N}\right)  \right)  \mathbf{,}\mathrm{w}_{\#}\right)
}{\delta\gamma}\varphi_{k}\left(  \mathbf{N}\right)  \right]  \left(
\mathbf{x}\right)  =\frac{\delta\mathcal{E}_{NI}\left(  \gamma\left(
\mathbf{\varphi}\left(  \mathbf{N}\right)  \right)  \mathbf{,}\mathrm{w}%
_{\#}\right)  }{\delta\gamma\left(  \mathbf{x}\right)  }\left[  \varphi
_{k}\left(  \mathbf{N}\right)  \right]  \left(  \mathbf{x}\right)  .
\end{equation}
The equation eq. $\left(  \ref{NIstat}\right)  ,$which could be expressed in
terms of the orthonormal GSO's $\mathbf{\tilde{\varphi}}\left(  \mathbf{N}%
\right)  $ as
\begin{equation}
\frac{\delta\mathcal{E}_{NI}\left(  \gamma\left(  \mathbf{\tilde{\varphi}%
}\left(  \mathbf{N}\right)  \right)  \mathbf{,}\mathrm{w}_{\#}\right)
}{\delta\gamma}\left\vert \tilde{\varphi}_{j}\left(  \mathbf{N}\right)
\right\rangle -\sum_{1\leq k\leq r_{1}}\tilde{\varepsilon}_{jk}\left(
\mathbf{N}\right)  \left\vert \tilde{\varphi}_{k}\left(  \mathbf{N}\right)
\right\rangle =0;\,1\leq j\leq r_{1},
\end{equation}
where
\begin{equation}
\tilde{\varepsilon}_{jk}\left(  \mathbf{N}\right)  =\varepsilon_{jk}\left(
\mathbf{N}\right)  \left(  \frac{N_{k}}{N_{j}}\right)  ^{\frac{1}{2}};\,1\leq
j,k\leq r_{1},
\end{equation}
is not an eigenvalue equation as $\mathcal{E}_{NI}\left(  \gamma\left(
\mathbf{\varphi}\left(  \mathbf{N}\right)  \right)  \mathbf{,}\mathrm{w}%
_{\#}\right)  $ is not invariant to mixing of GSO's that correspond to
different values of $N_{j}$'s. One also would expect that the minimum
$\mathcal{E}_{NI}\left(  \mathbf{\varphi}\left(  \mathbf{N}\right)
\mathbf{,}\mathrm{w}_{\#}\right)  $ will in general be degenerate i.e. many
values $\left\{  \mathbf{\varphi}_{\kappa}\left(  \mathbf{N}\right)  \right\}
$ will give this value. The isolated degenerate points are of course no
problem as they give different local minima.

The constraints on the control parameter $\mathbf{N}$
\begin{equation}
0\leq N_{j}\leq1,1\leq j\leq r_{1}%
\end{equation}
in the outer minimization of eq. $\left(  \ref{seqopt}\right)  $ can be held
by setting
\begin{equation}
N_{j}\left(  \mathbf{\theta}\right)  =\sin^{2}\theta_{j}%
\end{equation}
and introducing a Lagrange parameter $\eta$ for the constraint
\begin{equation}
\sum_{1\leq j\leq r_{1}}N_{j}\left(  \mathbf{\theta}\right)  =N.
\end{equation}
The outer control parameter $N$ could be \emph{any }postive number and
corresponds to the avarage number of electrons in the system. The constrained
minimum $E_{NIGS}\left(  \mathrm{w}_{\#}\right)  $ of $\mathcal{E}_{NI}\left(
\mathbf{\varphi}\left(  \mathbf{N}\right)  \mathbf{,}\mathrm{w}_{\#}\right)  $
with respect to $\mathbf{N}$ is then given by the saddle point of the
Lagrangian
\begin{equation}
\mathcal{L}_{N}\left(  \mathbf{N}\left(  \mathbf{\theta}\right)  ,\eta\left(
\mathbf{\theta}\right)  \mathbf{,}N\mathbf{,}\mathrm{w}_{\#}\right)
=\mathcal{E}_{NI}\left(  \mathbf{\varphi}\left(  \mathbf{N}\left(
\mathbf{\theta}\right)  \right)  \mathbf{,}\mathrm{w}_{\#}\right)
-\eta\left\{  \sum_{1\leq j\leq r_{1}}N_{j}-N\right\}  ,
\end{equation}
i.e.
\begin{equation}
\mathcal{L}_{N}\left(  \mathbf{N}\left(  \mathbf{\theta}\left(  N\right)
\right)  ,\eta\mathbf{,}\mathrm{w}_{\#}\right)  =\max_{\mathbf{\eta}}%
\min_{\mathbf{\theta}}\mathcal{L}_{N}\left(  \mathbf{\varphi}\left(
\mathbf{\theta}\right)  ,\eta\mathbf{,}N\mathbf{,}\mathrm{w}_{\#}\right)  .
\end{equation}
The necessary conditions for $\left(  \mathbf{\theta},\eta\right)  $ to be a
saddle point are
\begin{subequations}
\begin{align}
\frac{\partial\mathcal{L}_{N}\left(  \mathbf{\varphi}\left(  \mathbf{\theta
}\right)  ,\eta\mathbf{,}N\mathbf{,}\mathrm{w}_{\#}\right)  }{\partial
\theta_{j}}  &  =0,\,1\leq j\leq r_{1}\label{LNA}\\
\frac{\partial\mathcal{L}_{N}\left(  \mathbf{\varphi}\left(  \mathbf{\theta
}\right)  ,\eta\mathbf{,}N\mathbf{,}\mathrm{w}_{\#}\right)  }{\partial\eta}
&  =0. \label{LNB}%
\end{align}
Noting that
\end{subequations}
\begin{equation}
\frac{\partial\mathcal{L}_{N}\left(  \mathbf{\varphi}\left(  \mathbf{\theta
}\right)  ,\eta\mathbf{,}N\mathbf{,}\mathrm{w}_{\#}\right)  }{\partial
\theta_{j}}=\frac{\partial\mathcal{L}_{N}\left(  \mathbf{\varphi}\left(
\mathbf{\theta}\right)  ,\eta\mathbf{,}N\mathbf{,}\mathrm{w}_{\#}\right)
}{\partial N_{j}}\frac{\partial N_{j}}{\partial\theta_{j}}=\frac
{\partial\mathcal{L}_{N}\left(  \mathbf{\varphi}\left(  \mathbf{\theta
}\right)  ,\eta\mathbf{,}N\mathbf{,}\mathrm{w}_{\#}\right)  }{\partial N_{j}%
}\sin2\theta_{j};\,1\leq j\leq r_{1}%
\end{equation}
condition eq. $\left(  \ref{LNA}\right)  $ gives
\begin{equation}
\frac{\partial\mathcal{L}_{N}\left(  \mathbf{\varphi}\left(  \mathbf{\theta
}\right)  ,\eta\mathbf{,}N\mathbf{,}\mathrm{w}_{\#}\right)  }{\partial N_{j}%
}=0\text{ or }\sin2\theta_{j}=0.
\end{equation}
The derivative of the Lagrangian wrt $N_{j}$ produces%
\begin{equation}
\frac{\partial\mathcal{E}_{NI}\left(  \mathbf{\varphi}\left(  \mathbf{N}%
\left(  \mathbf{\theta}\right)  \right)  \mathbf{,}\mathrm{w}_{\#}\right)
}{\partial N_{j}}=\eta\left(  N\right)  ,
\end{equation}
which after observing
\begin{align}
&  \frac{\partial\mathcal{E}_{NI}\left(  \mathbf{\varphi}\left(
\mathbf{N}\left(  \mathbf{\theta}\right)  \right)  \mathbf{,}\mathrm{w}%
_{\#}\right)  }{\partial N_{j}}=\int\frac{\delta\mathcal{E}_{NI}\left(
\mathbf{\varphi}\left(  \mathbf{N}\left(  \mathbf{\theta}\right)  \right)
\mathbf{,}\mathrm{w}_{\#}\right)  }{\delta\gamma\left(  \mathbf{x}^{\prime
}\right)  }\frac{\partial\gamma\left(  \mathbf{x}^{\prime}\right)  }{\partial
N_{j}}d\mathbf{x}^{\prime}\nonumber\\
&  =\int\frac{\delta\mathcal{E}_{NI}\left(  \mathbf{\varphi}\left(
\mathbf{N}\left(  \mathbf{\theta}\right)  \right)  \mathbf{,}\mathrm{w}%
_{\#}\right)  }{\delta\gamma\left(  \mathbf{x}^{\prime}\right)  }\left[
\tilde{\varphi}_{j}\left(  \mathbf{N}\left(  \mathbf{\theta}\right)  \right)
\right]  \left(  \mathbf{x}^{\prime}\right)  ^{2}d\mathbf{x}^{\prime
}\nonumber\\
&  =\left\langle \tilde{\varphi}_{j}\left(  \mathbf{N}\left(  \mathbf{\theta
}\right)  \right)  \left\vert \frac{\delta\mathcal{E}_{NI}\left(
\mathbf{\varphi}\left(  \mathbf{N}\left(  \mathbf{\theta}\right)  \right)
\mathbf{,}\mathrm{w}_{\#}\right)  }{\delta\gamma}\tilde{\varphi}_{j}\left(
\mathbf{N}\left(  \mathbf{\theta}\right)  \right)  \right.  \right\rangle
=\tilde{\varepsilon}_{jj}\mathbf{\left(  \mathbf{N}\left(  \mathbf{\theta
}\right)  \right)  }=\varepsilon_{jj}\mathbf{\left(  \mathbf{N}\left(
\mathbf{\theta}\right)  \right)  },
\end{align}
one obtains
\begin{equation}
\frac{\partial\mathcal{E}_{NI}\left(  \mathbf{\varphi}\left(  \mathbf{N}%
\left(  \mathbf{\theta}\right)  \right)  \mathbf{,}\mathrm{w}_{\#}\right)
}{\partial N_{j}}=\varepsilon_{jj}\mathbf{\left(  \mathbf{N}\left(
\mathbf{\theta}\right)  \right)  }.
\end{equation}
After noting that
\begin{equation}
\sin2\theta_{j}=0\Rightarrow N_{j}\left(  \mathbf{\theta}\right)  =1,
\end{equation}
Hence
\begin{equation}
\frac{\partial\mathcal{E}_{NI}\left(  \mathbf{\varphi}\left(  \mathbf{N}%
\left(  \mathbf{\theta}\right)  \right)  \mathbf{,}\mathrm{w}_{\#}\right)
}{\partial N_{j}}=\varepsilon_{jj}\mathbf{\left(  \mathbf{N}\left(
\mathbf{\theta}\right)  \right)  }=\eta\left(  N\right)  \text{ or }%
N_{j}\left(  \mathbf{\theta}\right)  =1;\,1\leq j\leq r_{1}.
\end{equation}
If $N_{j}\left(  \mathbf{\theta}\right)  \neq1$ then $\frac{\partial
\mathcal{E}_{NI}\left(  \mathbf{\varphi}\left(  \mathbf{N}\left(
\mathbf{\theta}\right)  \right)  \mathbf{,}\mathrm{w}_{\#}\right)  }{\partial
N_{j}}=\varepsilon_{jj}\mathbf{\left(  \mathbf{N}\left(  \mathbf{\theta
}\right)  \right)  }=\eta\left(  N\right)  ,$ but if $\frac{\partial
\mathcal{E}_{NI}\left(  \mathbf{\varphi}\left(  \mathbf{N}\left(
\mathbf{\theta}\right)  \right)  \mathbf{,}\mathrm{w}_{\#}\right)  }{\partial
N_{j}}=\varepsilon_{jj}\mathbf{\left(  \mathbf{N}\left(  \mathbf{\theta
}\right)  \right)  }=\eta\left(  N\right)  $ then $N_{j}\left(  \mathbf{\theta
}\right)  $ might or might not equal one.

By considering sensitivity theory
\begin{equation}
\frac{\partial\mathcal{E}_{NI}\left(  \mathbf{\varphi}\left(  \mathbf{N}%
\left(  \mathbf{\theta}\right)  \right)  \mathbf{,}\mathrm{w}_{\#}\right)
}{\partial N}=\eta\left(  N\right)
\end{equation}
and $\eta\left(  N\right)  $ can be interpretted as the chemical potential of
the IP model system.

\subsubsection{Kohn-Sham energy functional}

One can define a Kohn-Sham energy functional, $\mathcal{F}_{KS},$ whose value
is the same as the HK energy functional, but whose decomposition is different
as
\begin{align}
\mathcal{F}_{KS}\left(  \gamma\right)   &  =\mathcal{F}_{KS}\left(
\Gamma\left(  \mathbf{\varphi}_{\#}\left(  \gamma,\mathbf{N}_{\#}\right)
\right)  \right)  =\mathcal{F}_{KS}\left(  \mathbf{\varphi}_{\#}\left(
\gamma,\mathbf{N}_{\#}\right)  \right) \nonumber\\
&  =\mathcal{F}_{T}\left(  \mathbf{\varphi}_{\#}\left(  \gamma,\mathbf{N}%
_{\#}\right)  \right)  +\mathcal{F}_{C}\left(  \gamma\right)  +\mathcal{F}%
_{KSXC}\left(  \mathbf{\varphi}_{\#}\left(  \gamma,\mathbf{N}_{\#}\right)
\right)  ,
\end{align}
where the Kohn-Sham exchange functional is defined as
\begin{equation}
\mathcal{F}_{KSXC}\left(  \mathbf{\varphi}_{\#}\left(  \gamma,\mathbf{N}%
_{\#}\right)  \right)  =\mathcal{F}_{XC}\left(  \gamma\right)  +\mathcal{F}%
_{K}\left(  \gamma\right)  -\mathcal{F}_{T}\left(  \mathbf{\varphi}%
_{\#}\left(  \gamma,\mathbf{N}_{\#}\right)  \right)  .
\end{equation}
(Recall $\mathcal{F}_{K}\left(  \gamma\right)  =\operatorname{Tr}\left\{
KD_{\ast}\left(  \gamma\right)  \right\}  $ and $\mathcal{F}_{XC}\left(
\gamma\right)  =\operatorname{Tr}\left\{  V_{2}D_{\ast}\left(  \gamma\right)
\right\}  -\mathcal{F}_{C}\left(  \gamma\right)  $)$.$The motivation for the
introduction of $\mathcal{F}_{KS}$ is that one can parameterize $\gamma$
uniquely in terms of $\left\{  \varphi_{\#j}\left(  \gamma,\mathbf{N}%
_{\#}\right)  \right\}  $ (which are unique up to unitary transformation that
mix orbitals that correspond to degenerate occupation numbers and occupation
numbers corresponding to degenerate density constrained ground states of $K$ )
through eq. $\left(  \ref{KEFunct}\right)  $ i.e.
\begin{equation}
\gamma\left(  \mathbf{x}\right)  =\sum_{1\leq j\leq r_{1}}\left\vert \left[
\varphi_{\#j}\left(  \gamma,\mathbf{N}_{\#}\right)  \right]  \left(
\mathbf{x}\right)  \right\vert ^{2}.
\end{equation}
If the correspondence was not unique then a different choice $\mathbf{\varphi
}\left(  \gamma,\mathbf{N}\right)  $ of orbitals would produce $\mathcal{F}%
_{KSXC}\left(  \mathbf{\varphi}\left(  \gamma,\mathbf{N}\right)  \right)
\neq\mathcal{F}_{KSXC}\left(  \mathbf{\varphi}_{\#}\left(  \gamma
,\mathbf{N}_{\#}\right)  \right)  $ as $\mathcal{F}_{T}\left(  \mathbf{\varphi
}\left(  \gamma,\mathbf{N}\right)  \right)  \neq\mathcal{F}_{T}\left(
\mathbf{\varphi}_{\#}\left(  \gamma,\mathbf{N}_{\#}\right)  \right)  .$In the
case of the degeneracy discussed before all the degenerate $\left\{
\mathbf{N}\right\}  $ produce the same values for Kohn-Sham exchange
functional so this does not result in such an ambiguity.

Actually the KS functional is based on a change of variable and a recipe for
the decomposition of the total energy functional. Unique values for KS XC
energy at any one density $\gamma$ are in fact \emph{not} necessary and could
vary depending on what set of GSO's is used to expand $\gamma$. A Kohn-Sham
energy functional can be expressed in terms of the variables $\left\{  \left.
\varphi_{j}\right\vert 1\leq j\leq r_{1}\right\}  $ as
\begin{equation}
\mathcal{F}_{KS}\left(  \mathbf{\varphi}\right)  =\mathcal{F}_{T}\left(
\mathbf{\varphi}\right)  +\mathcal{F}_{C}\left(  \mathbf{\varphi}\right)
+\mathcal{F}_{KSXC}\left(  \mathbf{\varphi}\right)  , \label{KS1}%
\end{equation}
where $\mathbf{\varphi}$ are sets of orthogonal functions belonging to the
unit ball of the \emph{real} Hilbert space $L_{2}\left(  \mathbb{X}\right)  ,$
and the density $\gamma=\Gamma\left(  \mathbf{\varphi}\right)  $ i.e. is
defined by the GSO's $\mathbf{\varphi}$ by
\begin{equation}
\gamma\left(  \mathbf{x}\right)  =\sum_{1\leq j\leq r_{1}}\varphi_{j}\left(
\mathbf{x}\right)  ^{2}. \label{DC}%
\end{equation}
The KS exchange corellation functional is defined as
\begin{equation}
\mathcal{F}_{KSXC}\left(  \mathbf{\varphi}\right)  =\mathcal{F}_{XC}\left(
\Gamma\left(  \mathbf{\varphi}\right)  \right)  +\mathcal{F}_{K}\left(
\mathbf{\varphi}\right)  -\mathcal{F}_{T}\left(  \Gamma\left(  \mathbf{\varphi
}\right)  \right)  , \label{KS2}%
\end{equation}
where,
\begin{equation}
\mathcal{F}_{XC}\left(  \Gamma\left(  \mathbf{\varphi}\right)  \right)
=\operatorname{Tr}\left\{  V_{2}D_{\ast}\left(  \gamma\right)  \right\}
-\mathcal{F}_{C}\left(  \gamma\right)  \label{KS3}%
\end{equation}
and
\begin{equation}
\mathcal{F}_{K}\left(  \mathbf{\varphi}\right)  =\operatorname{Tr}\left\{
KD_{\ast}\left(  \gamma\right)  \right\}  . \label{KS4}%
\end{equation}
One should note the functionals in eqs $\left(  \ref{KS3}\right)  $ and
$\left(  \ref{KS4}\right)  $ are constant for all $\mathbf{\varphi}$ that
satisfy $\gamma=\Gamma\left(  \mathbf{\varphi}\right)  $ for a fixed $\gamma$,
while the functional
\begin{equation}
\mathcal{F}_{T}\left(  \mathbf{\varphi}\right)  =\sum_{1\leq j\leq r_{1}%
}\left\langle \varphi_{j}\left\vert K\varphi_{j}\right.  \right\rangle
\end{equation}
can have different values even if $\Gamma\left(  \mathbf{\varphi}\right)
=\gamma$ and $\gamma$ is constant.

The minimization problem that determines the ground state energy and
space-spin density of a system characterized by a local potential $\mathrm{v}$
is
\begin{equation}
\min_{\mathbf{\varphi\in}\mathfrak{S}}\mathcal{E}_{KS}\left(  \mathbf{\varphi
,}\mathrm{v}\right)  \label{KSmin}%
\end{equation}
where
\begin{equation}
\mathcal{E}_{KS}\left(  \mathbf{\varphi,}\mathrm{v}\right)  =\mathcal{F}%
_{KS}\left(  \mathbf{\varphi}\right)  +f_{\Gamma\left(  \mathbf{\varphi
}\right)  }\left(  \mathrm{v}\right)  .
\end{equation}
The feasible set corresponds to the constraints
\begin{equation}
\left\langle \left.  \varphi_{j}\right\vert \varphi_{k}\right\rangle
=N_{j}\delta_{jk},\,1\leq j,k\leq r_{1},
\end{equation}
with control parameters $\left\{  \left.  N_{j}\right\vert 1\leq j,\leq
r_{1}\right\}  $ constrained by
\begin{equation}
0\leq N_{j}\leq1,\,1\leq j\leq r_{1};\sum_{\,1\leq j\leq r_{1}}N_{j}=N.
\end{equation}
The Kohn-Sham minimization problem can be decomposed into two sequential
optimizations, the first step is to determine the saddle point $\mathcal{L}%
\left(  \mathbf{\varphi}\left(  \mathbf{N}\right)  ,\mathbf{\varepsilon
}\left(  \mathbf{N}\right)  ,\mathrm{v}\right)  $ of the Largrangian
\begin{equation}
\mathcal{L}_{\mathbf{\varphi}}\left(  \mathbf{\varphi},\mathbf{\varepsilon
,N,}\mathrm{v}\right)  =\mathcal{E}_{KS}\left(  \mathbf{\varphi,}%
\mathrm{v}\right)  -\sum_{1\leq j,k\leq r_{1}}\varepsilon_{jk}\left\{
\overset{}{\left\langle \left.  \varphi_{j}\right\vert \varphi_{k}%
\right\rangle -N_{j}\delta_{jk}}\right\}  ,
\end{equation}
i.e.
\begin{equation}
\mathcal{L}_{\mathbf{\varphi}}\left(  \mathbf{\varphi}\left(  \mathbf{N}%
\right)  ,\mathbf{\varepsilon}\left(  \mathbf{N}\right)  ,\mathrm{v}\right)
=\max_{\mathbf{\varepsilon}}\min_{\mathbf{\varphi}}\mathcal{L}\left(
\mathbf{\varphi},\mathbf{\varepsilon,N,}\mathrm{v}\right)  , \label{KSM1}%
\end{equation}
leading to optimal values $\mathcal{E}_{KS}\left(  \mathbf{\varphi}\left(
\mathbf{N}\right)  \mathbf{,}\mathrm{v}\right)  =$ $\mathcal{L}%
_{\mathbf{\varphi}}\left(  \mathbf{\varphi}\left(  \mathbf{N}\right)
,\mathbf{\varepsilon}\left(  \mathbf{N}\right)  ,\mathbf{N,}\mathrm{v}\right)
$ of the KS energy for fixed values of the control parameters $\mathbf{N}$.
The second step to find a minima of $\mathcal{E}_{KS}\left(  \mathbf{\varphi
}\left(  \mathbf{N}\right)  \mathbf{,}\mathrm{v}\right)  $ with respect to a
constrained variation of these control parameters. The constraints
\begin{equation}
0\leq N_{j}\leq1,1\leq j\leq r_{1}%
\end{equation}
on the control parameter $\mathbf{N}$ can be held by setting
\begin{equation}
N_{j}\left(  \mathbf{\theta}\right)  =\sin^{2}\theta_{j},
\end{equation}
while one can introduce a Lagrange parameter $\eta$ in order to hold the
constraint
\begin{equation}
\sum_{1\leq j\leq r_{1}}N_{j}\left(  \mathbf{\theta}\right)  =N.
\end{equation}
This leads to a Lagrangian
\begin{equation}
\mathcal{L}_{\mathbf{\theta}}\left(  \mathbf{\varphi}\left(  \mathbf{N}\left(
\mathbf{\theta}\right)  \right)  ,\eta,N,\mathrm{v}\right)  =\mathcal{E}%
_{KS}\left(  \mathbf{\varphi}\left(  \mathbf{N}\right)  \mathbf{,}%
\mathrm{v}\right)  -\eta\left\{  \sum_{\,1\leq j\leq r_{1}}N_{j}-N\right\}
\end{equation}
and the determination of the saddle point $\mathcal{L}_{\mathbf{\theta}%
}\left(  \mathbf{\varphi}\left(  \mathbf{N}\left(  \mathbf{\theta}\left(
N\right)  \right)  \right)  ,\eta\left(  N\right)  ,\mathrm{v}\right)  $ given
by
\begin{equation}
\max_{\eta}\min_{\mathbf{\theta}}\mathcal{L}_{\mathbf{\theta}}\left(
\mathbf{\varphi}\left(  \mathbf{N}\left(  \mathbf{\theta}\right)  \right)
,\eta,N,\mathrm{v}\right)  .
\end{equation}
The saddle points of $\mathcal{L}_{\mathbf{\varphi}}\left(  \mathbf{\varphi
},\mathbf{\varepsilon,N,}\mathrm{v}\right)  $ are characterized by the
necessary conditions
\begin{subequations}
\begin{align}
\frac{\delta\mathcal{L}_{\mathbf{\varphi}}\left(  \mathbf{\varphi
},\mathbf{\varepsilon,N,}\mathrm{v}\right)  }{\delta\varphi_{j}}  &
=0;\,1\leq j\leq r_{1}\label{KSC1}\\
\frac{\partial\mathcal{L}_{\mathbf{\varphi}}\left(  \mathbf{\varphi
},\mathbf{\varepsilon,N,}\mathrm{v}\right)  }{\partial\varepsilon_{jk}}  &
=0;\,1\leq j,k\leq r_{1}. \label{KSC2}%
\end{align}
The condition eq. $\left(  \ref{KSC1}\right)  $ produces for each value of
$\mathbf{N}$
\end{subequations}
\begin{equation}
\frac{\delta\mathcal{E}_{KS}\left(  \mathbf{\varphi}\left(  \mathbf{N}\right)
,\mathbf{\varepsilon}\left(  \mathbf{N}\right)  ,\mathrm{v}\right)  }%
{\delta\varphi_{j}}-2\sum_{1\leq k\leq r_{1}}\varepsilon_{jk}\left(
\mathbf{N}\right)  \left\vert \varphi_{k}\left(  \mathbf{N}\right)
\right\rangle =0;\,1\leq j\leq r_{1}.
\end{equation}
which after noting that
\begin{align}
\frac{\delta\mathcal{E}_{KS}\left(  \mathbf{\varphi,}\mathrm{v}\right)
}{\delta\varphi_{j}\left(  \mathbf{x}\right)  }  &  =\int\frac{\delta
\mathcal{E}_{KS}\left(  \mathbf{\varphi,}\mathrm{v}\right)  }{\delta
\gamma\left(  \mathbf{x}^{\prime}\right)  }\frac{\delta\gamma\left(
\mathbf{x}^{\prime}\right)  }{\delta\varphi_{j}\left(  \mathbf{x}\right)
}d\mathbf{x}^{\prime}\nonumber\\
&  =2\int\frac{\delta\mathcal{E}_{KS}\left(  \mathbf{\varphi,}\mathrm{v}%
\right)  }{\delta\gamma\left(  \mathbf{x}\right)  }\varphi_{j}\left(
\mathbf{x}^{\prime}\right)  \delta\left(  \mathbf{x-x}^{\prime}\right)
d\mathbf{x}^{\prime}=2\frac{\delta\mathcal{E}_{KS}\left(  \mathbf{\varphi
,}\mathrm{v}\right)  }{\delta\gamma\left(  \mathbf{x}\right)  }\varphi
_{j}\left(  \mathbf{x}\right)
\end{align}
gives%
\begin{equation}
\frac{\delta\mathcal{E}_{KS}\left(  \mathbf{\mathbf{\varphi}\left(
\mathbf{N}\right)  ,}\mathrm{v}\right)  }{\delta\gamma}\left\vert \varphi
_{j}\left(  \mathbf{N}\right)  \right\rangle -\sum_{1\leq k\leq r_{1}%
}\varepsilon_{jk}\left(  \mathbf{N}\right)  \left\vert \varphi_{k}\left(
\mathbf{N}\right)  \right\rangle =0;\,1\leq j\leq r_{1}.
\end{equation}
The operator $\frac{\delta\mathcal{E}_{KS}\left(  \mathbf{\mathbf{\varphi
}\left(  \mathbf{N}\right)  ,}\mathrm{v}\right)  }{\delta\gamma}$ is a
multiplication operator given by
\begin{equation}
\frac{\delta\mathcal{E}_{KS}\left(  \mathbf{\mathbf{\varphi}\left(
\mathbf{N}\right)  ,}\mathrm{v}\right)  }{\delta\gamma}\left\vert \varphi
_{j}\left(  \mathbf{N}\right)  \right\rangle =\left\vert \frac{\delta
\mathcal{E}_{KS}\left(  \mathbf{\mathbf{\varphi}\left(  \mathbf{N}\right)
,}\mathrm{v}\right)  }{\delta\gamma}\varphi_{j}\left(  \mathbf{N}\right)
\right\rangle ,
\end{equation}
where
\begin{equation}
\left[  \frac{\delta\mathcal{E}_{KS}\left(  \mathbf{\mathbf{\varphi}\left(
\mathbf{N}\right)  ,}\mathrm{v}\right)  }{\delta\gamma}\varphi_{j}\left(
\mathbf{N}\right)  \right]  \left(  \mathbf{x}\right)  =\frac{\delta
\mathcal{E}_{NI}\left(  \gamma\left(  \mathbf{\varphi}\left(  \mathbf{N}%
\right)  \right)  \mathbf{,}\mathrm{w}_{\#}\right)  }{\delta\gamma\left(
\mathbf{x}\right)  }\left[  \varphi_{j}\left(  \mathbf{N}\right)  \right]
\left(  \mathbf{x}\right)  .
\end{equation}
We thus obtain the intermediate Kohn-Sham equations that correspond to
condition eq. $\left(  \ref{KSC1}\right)  $
\begin{equation}
\frac{\delta\mathcal{E}_{KS}\left(  \mathbf{\mathbf{\varphi}\left(
\mathbf{N}\right)  ,}\mathrm{v}\right)  }{\delta\gamma}\left\vert \varphi
_{j}\left(  \mathbf{N}\right)  \right\rangle =\sum_{1\leq k\leq r_{1}%
}\varepsilon_{jk}\left(  \mathbf{N}\right)  \left\vert \varphi_{k}\left(
\mathbf{N}\right)  \right\rangle .
\end{equation}
This is not an eigenvalue equation unless $\mathcal{E}_{KS}\left(
\mathbf{\varphi}\left(  \mathbf{N}\right)  ,\mathrm{v}\right)  $ is invariant
to unitary mixing of $\mathbf{\varphi}\left(  \mathbf{N}\right)  .$From
sensitivity theory the rate of change of the KS energy with respect to changes
in occupation numbers, which are the control parameters, is
\begin{equation}
\frac{\partial\mathcal{E}_{KS}\left(  \mathbf{\varphi}\left(  \mathbf{N}%
\right)  ,\mathrm{v}\right)  }{\partial N_{j}}=\varepsilon_{jj}\left(
\mathbf{N}\right)  .
\end{equation}
After noting that
\begin{align}
&  \frac{\partial\mathcal{L}_{\mathbf{\theta}}\left(  \mathbf{\varphi}\left(
\mathbf{N}\left(  \mathbf{\theta}\right)  \right)  ,\eta,N,\mathrm{v}\right)
}{\partial\theta_{j}}=\frac{\partial\mathcal{L}_{\mathbf{\theta}}\left(
\mathbf{\varphi}\left(  \mathbf{N}\left(  \mathbf{\theta}\right)  \right)
,\eta,N,\mathrm{v}\right)  }{\partial N_{j}}\frac{\partial N_{j}}%
{\partial\theta_{j}}=\frac{\partial\mathcal{L}_{\mathbf{\theta}}\left(
\mathbf{\varphi}\left(  \mathbf{N}\left(  \mathbf{\theta}\right)  \right)
,\eta,N,\mathrm{v}\right)  }{\partial N_{j}}\sin2\theta_{j}\,\nonumber\\
&  \qquad\qquad\qquad\qquad\qquad\qquad\qquad\qquad\qquad\qquad\qquad
\qquad\qquad\qquad\qquad1\leq j\leq r_{1},
\end{align}
one observes that the stationarity conditions eq. $\left(  \ref{KSC2}\right)
$ for the outer optimization give
\begin{equation}
\frac{\partial\mathcal{L}_{\mathbf{\theta}}\left(  \mathbf{\varphi}\left(
\mathbf{N}\left(  \mathbf{\theta}\left(  N\right)  \right)  \right)
,\eta\mathbf{,}N\mathbf{,}\mathrm{v}\right)  }{\partial N_{j}}=0\text{ }
\label{out1}%
\end{equation}
or%
\begin{equation}
\sin2\theta\left(  N\right)  _{j}=0. \label{out2}%
\end{equation}
The stationary condition eq. $\left(  \ref{out1}\right)  $ produces
\begin{equation}
\frac{\partial\mathcal{E}_{KS}\left(  \mathbf{\varphi}\left(  \mathbf{N}%
\left(  \mathbf{\theta}\left(  N\right)  \right)  \right)  \mathbf{,}%
\mathrm{v}\right)  }{\partial N_{j}}=\eta\left(  N\right)  ,
\end{equation}
which after noting
\begin{align}
&  \frac{\partial\mathcal{E}_{KS}\left(  \mathbf{\varphi}\left(
\mathbf{N}\right)  \mathbf{,}\mathrm{v}\right)  }{\partial N_{j}}=\int
\frac{\delta\mathcal{E}_{KS}\left(  \mathbf{\varphi}\left(  \mathbf{N}\right)
\mathbf{,}\mathrm{v}\right)  }{\delta\gamma\left(  \mathbf{x}^{\prime}\right)
}\frac{\partial\gamma\left(  \mathbf{x}^{\prime}\right)  }{\partial N_{j}%
}d\mathbf{x}^{\prime}\nonumber\\
&  =\int\frac{\delta\mathcal{E}_{KS}\left(  \mathbf{\varphi}\left(
\mathbf{N}\right)  \mathbf{,}\mathrm{v}\right)  }{\delta\gamma\left(
\mathbf{x}^{\prime}\right)  }\left[  \tilde{\varphi}_{j}\left(  \mathbf{N}%
\left(  \mathbf{\theta}\right)  \right)  \right]  \left(  \mathbf{x}^{\prime
}\right)  ^{2}d\mathbf{x}^{\prime}\nonumber\\
&  =\left\langle \tilde{\varphi}_{j}\left(  \mathbf{N}\right)  \left\vert
\frac{\delta\mathcal{E}_{KS}\left(  \mathbf{\varphi}\left(  \mathbf{N}\right)
\mathbf{,}\mathrm{v}\right)  }{\delta\gamma}\tilde{\varphi}_{j}\left(
\mathbf{N}\right)  \right.  \right\rangle =\tilde{\varepsilon}_{jj}%
\mathbf{\left(  \mathbf{N}\right)  }=\varepsilon_{jj}\mathbf{\left(
\mathbf{N}\right)  },
\end{align}
one obtains
\begin{equation}
\frac{\partial\mathcal{E}_{KS}\left(  \mathbf{\varphi}\left(  \mathbf{N}%
\left(  \mathbf{\theta}\left(  N\right)  \right)  \right)  \mathbf{,}%
\mathrm{v}\right)  }{\partial N_{j}}=\varepsilon_{jj}\mathbf{\left(
\mathbf{N}\left(  \mathbf{\theta}\left(  N\right)  \right)  \right)  }%
=\eta\left(  N\right)  .
\end{equation}
The stationary condition eq. $\left(  \ref{out2}\right)  $ gives
\begin{equation}
\sin2\theta_{j}\left(  N\right)  =0\Rightarrow N_{j}\left(  \mathbf{\theta
}\left(  N\right)  \right)  =1.
\end{equation}
Hence
\begin{equation}
\frac{\partial\mathcal{E}_{KS}\left(  \mathbf{\varphi}\left(  \mathbf{N}%
\left(  \mathbf{\theta}\left(  N\right)  \right)  \right)  \mathbf{,}%
\mathrm{v}\right)  }{\partial N_{j}}=\eta\left(  N\right)  \text{ or }%
N_{j}\left(  \mathbf{\theta}\left(  N\right)  \right)  =1;\,1\leq j\leq r_{1}.
\end{equation}
If $N_{j}\left(  \mathbf{\theta}\left(  N\right)  \right)  \neq1$ then
$\frac{\partial\mathcal{E}_{KS}\left(  \mathbf{\varphi}\left(  \mathbf{N}%
\left(  \mathbf{\theta}\left(  N\right)  \right)  \right)  \mathbf{,}%
\mathrm{v}\right)  }{\partial N_{j}}=\varepsilon_{jj}\mathbf{\left(
\mathbf{N}\left(  \mathbf{\theta}\left(  N\right)  \right)  \right)  }%
=\eta\left(  N\right)  ,$ but if $\frac{\partial\mathcal{E}_{KS}\left(
\mathbf{\varphi}\left(  \mathbf{N}\left(  \mathbf{\theta}\left(  N\right)
\right)  \right)  \mathbf{,}\mathrm{v}\right)  }{\partial N_{j}}%
=\varepsilon_{jj}\mathbf{\left(  \mathbf{N}\left(  \mathbf{\theta}\left(
N\right)  \right)  \right)  }=\eta\left(  N\right)  $ then $N_{j}\left(
\mathbf{\theta}\left(  N\right)  \right)  $ might or might not equal one.

By considering sensitivity theory one obtains
\begin{equation}
\frac{\partial\mathcal{E}_{KS}\left(  \mathbf{\varphi}\left(  \mathbf{N}%
\left(  \mathbf{\theta}\left(  N\right)  \right)  \right)  \mathbf{,}%
\mathrm{v}\right)  }{\partial N}=\eta\left(  N\right)
\end{equation}
and $\eta\left(  N\right)  $ can be interpretted as the chemical potential of
the interacting molecular system.

\subsubsection{Janaks Theorem}

One can develop the Kohn-Sham energy functional as%
\begin{align}
&  \mathcal{E}_{KS}\left(  \mathbf{\varphi}\left(  \mathbf{N}\right)
\mathbf{,}\mathrm{v}\right)  =\mathcal{F}_{KS}\left(  \mathbf{\varphi}\left(
\mathbf{N}\right)  \right)  +f_{\Gamma\left(  \mathbf{\varphi}\right)
}\left(  \mathrm{v}\right) \nonumber\\
&  =\mathcal{F}_{T}\left(  \mathbf{\varphi}\left(  \mathbf{N}\right)  \right)
+\mathcal{F}_{C}\left(  \mathbf{\varphi}\left(  \mathbf{N}\right)  \right)
+\mathcal{F}_{KSXC}\left(  \mathbf{\varphi}\left(  \mathbf{N}\right)  \right)
+f_{\Gamma\left(  \mathbf{\varphi}\right)  }\left(  \mathrm{v}\right)
\nonumber\\
&  =\sum_{1\leq j\leq r_{1}}\left\langle \widetilde{\varphi_{j}}\left\vert
K\widetilde{\varphi_{j}}\right.  \right\rangle N_{j}+\int\frac{\gamma\left(
\mathbf{r,\xi}\right)  \gamma\left(  \mathbf{r}^{\prime}\mathbf{,\xi}\right)
}{\left\Vert \mathbf{r-r}^{\prime}\right\Vert }d\mathbf{r}d\mathbf{r}^{\prime
}d\mathbf{\xi}+\int\mathrm{v}\left(  \mathbf{r,\xi}\right)  \gamma\left(
\mathbf{r,\xi}\right)  d\mathbf{r}d\mathbf{\xi}+\mathcal{F}_{KSXC}\left(
\mathbf{\varphi}\left(  \mathbf{N}\right)  \right)  ,
\end{align}
where
\begin{equation}
\gamma\left(  \mathbf{x}\right)  =\sum_{1\leq j\leq r_{1}}\varphi_{j}\left(
\mathbf{x}\right)  ^{2}=\sum_{1\leq j\leq r_{1}}\widetilde{\varphi_{j}\left(
\mathbf{x}\right)  }^{2}N_{j}.
\end{equation}
Using
\begin{equation}
\frac{\partial\mathcal{E}_{KS}\left(  \mathbf{\varphi}\left(  \mathbf{N}%
\right)  ,\mathrm{v}\right)  }{\partial N_{j}}=\int\frac{\delta\mathcal{E}%
_{KS}\left(  \mathbf{\varphi}\left(  \mathbf{N}\right)  ,\mathrm{v}\right)
}{\delta\gamma\left(  \mathbf{x}\right)  }\frac{\partial\gamma\left(
\mathbf{x}\right)  }{\partial N_{j}}d\mathbf{x=}\int\frac{\delta
\mathcal{E}_{KS}\left(  \mathbf{\varphi}\left(  \mathbf{N}\right)
,\mathrm{v}\right)  }{\delta\gamma\left(  \mathbf{x}\right)  }\widetilde
{\varphi_{j}\left(  \mathbf{x}\right)  }^{2}d\mathbf{x} \label{DerNj}%
\end{equation}
one obtains Janaks Theorem
\begin{align}
&  \frac{\partial\mathcal{E}_{KS}\left(  \mathbf{\varphi}\left(
\mathbf{N}\right)  ,\mathrm{v}\right)  }{\partial N_{j}}=\left\langle
\widetilde{\varphi_{j}}\left\vert \frac{\delta\mathcal{E}_{KS}\left(
\mathbf{\varphi}\left(  \mathbf{N}\right)  ,\mathrm{v}\right)  }{\delta
\gamma\left(  \mathbf{x}\right)  }\widetilde{\varphi_{j}}\right.
\right\rangle =\nonumber\\
&  \left\langle \widetilde{\varphi_{j}}\left\vert K\widetilde{\varphi_{j}%
}\right.  \right\rangle +\int\frac{\delta}{\delta\gamma\left(  \mathbf{x}%
\right)  }\left\{  \mathcal{F}_{C}\left(  \mathbf{\varphi}\left(
\mathbf{N}\right)  \right)  +\mathcal{F}_{KSXC}\left(  \mathbf{\varphi}\left(
\mathbf{N}\right)  \right)  +f_{\Gamma\left(  \mathbf{\varphi}\right)
}\left(  \mathrm{v}\right)  \right\}  \frac{\partial\gamma\left(
\mathbf{x}\right)  }{\partial N_{j}}d\mathbf{x=}\nonumber\\
&  \left\langle \widetilde{\varphi_{j}}\left\vert K\widetilde{\varphi_{j}%
}\right.  \right\rangle +\left\langle \widetilde{\varphi_{j}}\left\vert
\frac{\delta}{\delta\gamma\left(  \mathbf{x}\right)  }\left\{  \mathcal{F}%
_{C}\left(  \mathbf{\varphi}\left(  \mathbf{N}\right)  \right)  +\mathcal{F}%
_{KSXC}\left(  \mathbf{\varphi}\left(  \mathbf{N}\right)  \right)
+f_{\Gamma\left(  \mathbf{\varphi}\right)  }\left(  \mathrm{v}\right)
\right\}  \right.  \widetilde{\varphi_{j}}\right\rangle =\varepsilon
_{jj}\left(  \mathbf{N}\right)  .
\end{align}


\section{AGP Energy Functionals\label{AGP}}

\subsection{Canonical expansion of an AGP state}

\bigskip The $2N$-electron AGP state, $\left\vert g^{N}\right\rangle $, is the
$N^{th}$ antisymmetric power of the $2$-electron state described by the
geminal
\begin{equation}
\left\vert g\right\rangle =\sum_{1\leq j\leq s}n_{j}^{\frac{1}{2}}\left\vert
\psi_{j}\psi_{j+s}\right\rangle , \label{geminal}%
\end{equation}
where $\left\{  \left.  n_{j}^{\frac{1}{2}}\right\vert 1\leq j\leq s\right\}
$ are the positive square roots of $\left\{  \left.  n_{j}\right\vert 1\leq
j\leq s\right\}  ,$ and is given by \cite{Weiner&ortiz04}
\begin{equation}
\left\vert g^{N}\right\rangle =\left(  S_{N}\right)  ^{-\frac{1}{2}}%
\sum_{1\leq j_{1}<\cdots<j_{N}\leq s}n_{j_{1}}^{\frac{1}{2}}\cdots n_{j_{s}%
}^{\frac{1}{2}}\left\vert \psi_{j_{1}}\psi_{j_{1}+s}\cdots\psi_{j_{N}}%
\psi_{j_{N}+s}\right\rangle . \label{agp}%
\end{equation}
The normalization factor, $\left(  S_{N}\right)  ^{-\frac{1}{2}}$, is given in
terms of the elementary symmetric function
\begin{equation}
S_{N}=\sum_{1\leq j_{1}<\cdots<j_{N}\leq s}n_{j_{1}}\cdots n_{j_{N}},
\end{equation}
and the geminals $\left\vert g\right\rangle $ represent normalized
two--electron states so that
\begin{equation}
\sum_{1\leq j\leq s}n_{j}=1.
\end{equation}
Note that
\begin{equation}
S_{0}=1.
\end{equation}
It is evident from eq.\ $\left(  \ref{geminal}\right)  $ that the geminal and
geminal power state are at the least invariant to the group
\begin{equation}
SU\left(  2\right)  ^{s}=\overset{s\text{-factors}}{\overbrace{\,SU\left(
2\right)  \times\cdots\times SU\left(  2\right)  }},
\end{equation}
that is formed by $2\times2$ \emph{special} unitary transformations of the
paired CGSOs $\left\{  \left.  \left\vert \psi_{j}\right\rangle ,\left\vert
\psi_{j+s}\right\rangle \right\vert 1\leq j\leq s\right\}  ,$ while if
degeneracies exist amongst the $\left\{  \left.  n_{j}\right\vert 1\leq j\leq
s\right\}  $ this invariance group is larger \cite{AGPCoherent}. The canonical
CSGOs are unique up to transformations belonging this invariance group.

\subsection{The AGP Energy Functional}

The energy of the AGP state is given by \emph{ }
\begin{align}
\mathcal{E}\left(  \mathbf{n,\varphi}\right)   &  =\frac{\mathcal{H}\left(
\mathbf{n,\varphi}\right)  }{S_{N}\left(  \mathbf{n}\right)  }\nonumber\\
&  =\frac{1}{S_{N}\left(  \mathbf{n}\right)  }\left\{  \sum_{1\leq j\leq
s}n_{j}S_{N-1}\left(  \jmath\right)  \left\{  \left\langle \varphi
_{j}\left\vert h_{1}\varphi_{j}\right.  \right\rangle +\left\langle
\varphi_{j+s}\left\vert h_{1}\varphi_{j+s}\right.  \right\rangle +\left\langle
\varphi_{j}\varphi_{j+s}||\varphi_{j}\varphi_{j+s}\right\rangle \right\}
\right. \nonumber\\
&  +\sum_{1\leq j,k\leq s}n_{j}^{\frac{1}{2}}n_{k}^{\frac{1}{2}}S_{N-1}\left(
\jmath k\right)  \left\langle \varphi_{j}\varphi_{j+s}||\varphi_{k}%
\varphi_{k+s}\right\rangle \left.  +\sum_{1\leq j<k\leq s}n_{i}n_{j}%
S_{N-2}\left(  \jmath k\right)  \left\langle \varphi_{j}\varphi_{k}%
|||\varphi_{j}\varphi_{k}\right\rangle \right\}  \label{energy}%
\end{align}
in terms of geminal occupation numbers and orthonormal CGSO's, while in terms
of the one electron group parameters it is
\begin{equation}
\mathcal{E}\left(  \mathbf{\eta},\mathbf{\lambda}\right)  =\frac
{\mathcal{H}\left(  \mathbf{\eta},\mathbf{\lambda}\right)  }{S_{N}\left(
\mathbf{\eta}\right)  }=\frac{\left\langle \left.  g\left(  \mathbf{\eta
},\mathbf{\lambda}\right)  ^{N}\right\vert Hg\left(  \mathbf{\eta
},\mathbf{\lambda}\right)  ^{N}\right\rangle }{S_{N}\left(  \mathbf{\eta
}\right)  }, \label{gp_funct}%
\end{equation}
where%
\begin{equation}
S_{N}\left(  \mathbf{\eta}\right)  =\left\langle \left.  k\left(
\mathbf{\eta}\right)  g_{0}^{N}\right\vert k\left(  \mathbf{\eta}\right)
g_{0}^{N}\right\rangle =\sum_{1\leq j_{1}<\cdots<j_{N}\leq s}e^{2\left(
\eta_{j_{1}}+\cdots+\eta_{j_{N}}\right)  }n_{0j_{1}}\cdots n_{0j_{N}}.
\end{equation}


\section{\bigskip AGP Lagrangian\label{Lagran}}

One can define a Lagrangian
\begin{equation}
\mathcal{L}\left(  \mathbf{n,\varphi}\right)  =\mathcal{E}\left(
\mathbf{n,\varphi}\right)  -\sum_{1\leq j,k\leq2s}\varepsilon_{kj}\left(
\left\langle \left.  \varphi_{j}\right\vert \varphi_{k}\right\rangle
-\delta_{jk}\right)
\end{equation}
associated with the constrained variation (orthogonal with a fixed
normalization) of energy function $\mathcal{E}\left(  \mathbf{n,\varphi
}\right)  .$ At constrained stationary points
\begin{equation}
\frac{\partial\mathcal{L}\left(  \mathbf{n,\varphi}\right)  }{\partial n_{j}%
}=\frac{\partial\mathcal{E}\left(  \mathbf{n,\varphi}\right)  }{\partial
n_{j}}.
\end{equation}
This gives
\begin{align}
&  \frac{\partial\mathcal{L}\left(  \mathbf{n,\varphi}\right)  }{\partial
N_{j}}=\sum_{1\leq k\leq2s}\frac{\partial\mathcal{L}\left(  \mathbf{n,\varphi
}\right)  }{\partial n_{k}}\frac{\partial n_{k}}{\partial N_{j}}\\
&  \Rightarrow\frac{\partial\mathcal{L}\left(  \mathbf{n,\varphi}\right)
}{\partial N_{j}}=\sum_{1\leq k\leq2s}\frac{\partial\mathcal{L}\left(
\mathbf{n,\varphi}\right)  }{\partial n_{k}}\frac{\partial n_{k}}{\partial
N_{j}}=\sum_{1\leq k\leq2s}\frac{\partial\mathcal{L}\left(  \mathbf{n,\varphi
}\right)  }{\partial n_{k}}\left(  \frac{\partial N_{k}}{\partial n_{j}%
}\right)  ^{-1}%
\end{align}
i.e.%
\begin{equation}
\frac{\partial\mathcal{L}\left(  \mathbf{n,\varphi}\right)  }{\partial
\mathbf{N}}=\frac{\partial\mathcal{L}\left(  \mathbf{n,\varphi}\right)
}{\partial\mathbf{N}}=\frac{\partial\mathcal{L}\left(  \mathbf{n,\varphi
}\right)  }{\partial\mathbf{n}}\left(  \frac{\partial\mathbf{N}}%
{\partial\mathbf{n}}\right)  ^{-1}.
\end{equation}
At a constrained stationary point
\begin{equation}
\frac{\partial\mathcal{L}\left(  \mathbf{n,\varphi}\right)  }{\partial
\mathbf{N}}=\frac{\partial\mathcal{L}\left(  \mathbf{n,\varphi}\right)
}{\partial\mathbf{n}}=\mathbf{0.}%
\end{equation}
The occupation numbers of the AGP\ state are
\begin{align}
N_{j}  &  =\mathcal{N}_{j}\left(  \mathbf{n}\right)  =n_{j}\frac
{S_{N-1}\left(  \jmath\right)  }{S_{N}}=n_{j}\frac{\partial\ln S_{N}}{\partial
n_{j}};\quad1\leq j\leq s,\\
\mathbf{N}  &  =\mathcal{N}\left(  \mathbf{n}\right)  .
\end{align}%
\begin{equation}
\frac{\partial S_{N}}{\partial n_{j}}=S_{N-1}\left(  j\right)
\end{equation}%
\begin{align}
\sum_{j}n_{j}S_{N-1}\left(  j\right)   &  =NS_{N}\\
\sum_{k}n_{k}S_{N-2}\left(  jk\right)   &  =\left(  N-1\right)  S_{N-1}\left(
j\right) \\
\sum_{j<k}n_{j}n_{k}S_{N-2}\left(  jk\right)   &  =\frac{N\left(  N-1\right)
}{2}S_{N}%
\end{align}
The Jacobian is given by
\begin{equation}
\frac{\partial\mathcal{L}\left(  \mathbf{n,\varphi}\right)  }{\partial n_{j}%
}=\sum_{k}\frac{\partial\mathcal{L}\left(  \mathbf{n,\varphi}\right)
}{\partial N_{k}}\frac{\partial N_{k}}{\partial n_{j}}=\sum_{k}\frac
{\partial\mathcal{L}\left(  \mathbf{n,\varphi}\right)  }{\partial N_{k}}J_{kj}%
\end{equation}
The matrix elements of the Jacobian are
\begin{align}
\frac{\partial N_{k}}{\partial n_{j}}  &  =\frac{\partial}{\partial n_{j}%
}\left\{  n_{k}\frac{\partial\ln S_{N}}{\partial n_{k}}\right\}  =\delta
_{jk}\frac{\partial\ln S_{N}}{\partial n_{k}}+n_{k}\frac{\partial^{2}\ln
S_{N}}{\partial n_{j}\partial n_{k}}\\
&  =\delta_{jk}\frac{S_{N-1}\left(  \jmath\right)  }{S_{N}}-n_{k}\frac
{S_{N-1}\left(  \jmath\right)  S_{N-1}\left(  k\right)  }{S_{N}^{2}}%
+n_{k}\frac{S_{N-2}\left(  jk\right)  }{S_{N}}.\nonumber
\end{align}
For a $2N$-electron IPS the geminal has the form
\begin{equation}
\left\vert g\right\rangle =N^{-\frac{1}{2}}\sum_{1\leq j\leq N}\left\vert
\psi_{j}\psi_{j+s}\right\rangle
\end{equation}
which for a $8$-electron case is
\begin{equation}
\left\vert g\right\rangle =4^{-\frac{1}{2}}\sum_{1\leq j\leq4}\left\vert
\psi_{j}\psi_{j+s}\right\rangle =\frac{1}{2}\sum_{1\leq j\leq4}\left\vert
\psi_{j}\psi_{j+s}\right\rangle
\end{equation}
so that
\begin{equation}
n_{j}=\left\{
\begin{array}
[c]{c}%
N^{-1};1\leq j\leq N\\
\qquad0;\text{ }N+1\leq j\leq s
\end{array}
\right.
\end{equation}
which for a $6$-electron case is%
\begin{equation}
n_{j}=\left\{
\begin{array}
[c]{c}%
\frac{1}{4};1\leq j\leq4\\
0;\text{ }5\leq j\leq s
\end{array}
\right.
\end{equation}%
\begin{equation}
S_{N}=N^{-N}%
\end{equation}
which for a $8$-electron case is%
\begin{equation}
S_{4}=\frac{1}{256}.
\end{equation}
The first derivatives of the symmetric function are
\begin{equation}
S_{N-1}\left(  \jmath\right)  =\sum_{\substack{j\notin\left\{  j_{1}%
,\cdots,j_{N-1}\right\}  \\1\leq j_{1}<\cdots<j_{N-1}\leq N}}n_{j_{1}}\cdots
n_{j_{N-1}}=\left\{
\begin{array}
[c]{c}%
\binom{N}{N-1}N^{-\left(  N-1\right)  };\quad j\notin\left\{  1,\cdots
,N\right\} \\
\binom{N-1}{N-1}N^{-\left(  N-1\right)  };\quad j\in\left\{  1,\cdots
,N\right\}
\end{array}
\right.
\end{equation}
which for a $8$-electron $s=8$ case where $\left(  n_{1},n_{2},n_{3}%
,n_{4}\right)  =\left(  \frac{1}{4},\frac{1}{4},\frac{1}{4},\frac{1}%
{4}\right)  $ is%

\begin{align}
S_{3}\left(  8\right)   &  =n_{1}n_{2}n_{3}+n_{1}n_{2}n_{4}+n_{1}n_{2}%
n_{5}+n_{1}n_{2}n_{6}+n_{1}n_{2}n_{7}\\
&  +n_{1}n_{3}n_{4}+n_{1}n_{3}n_{5}+n_{1}n_{3}n_{6}+n_{1}n_{3}n_{7}\\
&  +n_{1}n_{4}n_{5}+n_{1}n_{4}n_{6}+n_{1}n_{4}n_{7}\\
&  +n_{1}n_{5}n_{6}+n_{1}n_{5}n_{7}+n_{1}n_{6}n_{7}\\
&  +n_{2}n_{3}n_{4}+n_{2}n_{3}n_{5}+n_{2}n_{3}n_{6}+n_{2}n_{3}n_{7}\\
&  +n_{2}n_{4}n_{5}+n_{2}n_{4}n_{6}+n_{2}n_{4}n_{7}\\
&  +n_{2}n_{5}n_{6}+n_{2}n_{5}n_{7}+n_{2}n_{6}n_{7}\\
&  +n_{3}n_{4}n_{5}+n_{3}n_{4}n_{6}+n_{3}n_{4}n_{7}\\
&  +n_{3}n_{5}n_{6}+n_{3}n_{5}n_{7}+n_{3}n_{6}n_{7}\\
&  +n_{4}n_{5}n_{6}+n_{4}n_{5}n_{7}+n_{4}n_{6}n_{7}+n_{5}n_{6}n_{7}%
\end{align}%
\begin{equation}
S_{4}\left(  \jmath\right)  =\sum_{\substack{j\notin\left\{  j_{1},j_{2}%
,j_{3},j_{4}\right\}  \\1\leq j_{1}<j_{2}<j_{3}<j_{4}\leq8}}n_{j_{1}}n_{j_{2}%
}n_{j_{3}}n_{j_{4}}=\left\{
\begin{array}
[c]{c}%
\frac{1}{16};\quad j\notin\left\{  1,2,3,4\right\} \\
\frac{1}{64};\quad j\in\left\{  1,2,3,4\right\}
\end{array}
\right.
\end{equation}%
\begin{equation}
S_{N-2}\left(  jk\right)  =\sum_{\substack{j,k\notin\left\{  j_{1}%
,\cdots,j_{N-2}\right\}  \\1\leq j_{1}<\cdots<j_{N-2}\leq N}}n_{j_{1}}\cdots
n_{j_{N-2}}=\left\{
\begin{array}
[c]{c}%
\binom{N}{N-2}N^{-\left(  N-2\right)  };\quad j,k\notin\left\{  1,\cdots
,N\right\} \\
\binom{N-1}{N-2}N^{-\left(  N-2\right)  };\quad j\notin\left\{  1,\cdots
,N\right\}  ,k\in\left\{  1,\cdots,N\right\} \\
\binom{N-1}{N-2}N^{-\left(  N-2\right)  };\quad k\notin\left\{  1,\cdots
,N\right\}  ,j\in\left\{  1,\cdots,N\right\} \\
\binom{N-2}{N-2}N^{-\left(  N-2\right)  };\quad j\in\left\{  1,\cdots
,N\right\}  ,k\in\left\{  1,\cdots,N\right\}
\end{array}
\right.
\end{equation}
which for a $8$-electron case is
\begin{align}
S_{2}\left(  7,8\right)  =n_{1}n_{2}  &  +n_{1}n_{3}+n_{1}n_{4}+n_{1}%
n_{5}+n_{1}n_{6}+n_{2}n_{3}+n_{2}n_{4}+n_{2}n_{5}\\
&  +n_{2}n_{6}+n_{3}n_{4}+n_{3}n_{5}+n_{3}n_{6}+n_{4}n_{5}+n_{4}n_{6}%
+n_{5}n_{6}%
\end{align}%
\begin{equation}
S_{2}\left(  jk\right)  =\sum_{\substack{j,k\notin\left\{  j_{1}%
,\cdots,j_{N-2}\right\}  \\1\leq j_{1}<\cdots<j_{N-2}\leq N}}n_{j_{1}}\cdots
n_{j_{N-2}}=\left\{
\begin{array}
[c]{c}%
\frac{3}{8};\quad j,k\notin\left\{  1,\cdots,N\right\}  \qquad\qquad\qquad\\
\frac{3}{16};\quad j\notin\left\{  1,\cdots,N\right\}  ,k\in\left\{
1,\cdots,N\right\} \\
\frac{3}{16};\quad k\notin\left\{  1,\cdots,N\right\}  ,j\in\left\{
1,\cdots,N\right\} \\
\frac{1}{16};\quad j\in\left\{  1,\cdots,N\right\}  ,k\in\left\{
1,\cdots,N\right\}
\end{array}
\right.
\end{equation}
The diagonal
\begin{equation}
\frac{\partial N_{k}}{\partial n_{k}}=\left\{
\begin{array}
[c]{c}%
N;\quad k\notin\left\{  1,\cdots,N\right\} \\
0;\quad k\in\left\{  1,\cdots,N\right\}
\end{array}
\right.
\end{equation}%
\begin{equation}
\frac{\partial N_{k}}{\partial n_{j}}=\left\{
\begin{array}
[c]{c}%
-n_{k}\frac{N\left(  N+1\right)  }{2};\quad j\neq k\notin\left\{
1,\cdots,N\right\} \\
n_{k};\quad j\notin\left\{  1,\cdots,N\right\}  ,k\in\left\{  1,\cdots
,N\right\} \\
n_{k};\quad k\notin\left\{  1,\cdots,N\right\}  ,j\in\left\{  1,\cdots
,N\right\} \\
0;\quad j\in\left\{  1,\cdots,N\right\}  ,k\in\left\{  1,\cdots,N\right\}
\end{array}
\right.
\end{equation}


The occupation number variation nonunitary subgroup $G_{num}$ is given by
\begin{equation}
G_{num}=\left\{  \left.  k\left(  \mathbf{\eta}\right)  =\exp\left\{
\sum_{1\leq j\leq s}\eta_{j}\Omega_{j}\right\}  \right\vert \eta_{j}%
\in\mathbb{R}\right\}  , \label{occup}%
\end{equation}
where
\begin{equation}
\Omega_{j}=\left(  a_{j}^{\dagger}a_{j}+a_{j+s}^{\dagger}a_{j+s}\right)
;\quad1\leq j\leq s.
\end{equation}
The manifold of $2N$-electron AGP\ states $\mathcal{M}$ can be expressed as
\begin{equation}
\mathcal{M}\mathcal{=}\left\{  \left.  \overset{}{k\left(  \mathbf{\eta
}\right)  u\left(  \mathbf{\lambda}\right)  }\left\vert g_{0}^{N}\right\rangle
\right\vert \right.  \left.  \mathbf{\lambda}=\mathbf{\lambda}^{\dagger
}\right\}  .
\end{equation}
Explicitly the action on a reference geminal is given by
\begin{equation}
\left\vert \tilde{g}\left(  \mathbf{\eta},\mathbf{\lambda}\right)
\right\rangle =k\left(  \mathbf{\eta}\right)  u\left(  \mathbf{\lambda
}\right)  \left\vert g_{0}\right\rangle =\sum_{1\leq l\leq s}e^{\eta_{l}%
}c_{0l}\left\vert \varphi_{l}\left(  \mathbf{\lambda}\right)  \varphi
_{l+s}\left(  \mathbf{\lambda}\right)  \right\rangle ,
\end{equation}
where
\begin{equation}
\left\vert \varphi_{l}\left(  \mathbf{\lambda}\right)  \right\rangle
=\exp\left\{  i\sum_{1\leq j,k\leq2s}\lambda_{jk}a_{j}^{\dagger}a_{k}\right\}
\left\vert \varphi_{0l}\right\rangle .
\end{equation}
One can define a normalized geminal by
\begin{equation}
\left\vert g\left(  \mathbf{\eta},\mathbf{\lambda}\right)  \right\rangle
=\mathit{s}\left(  \mathbf{\eta}\right)  ^{-\frac{1}{2}}\left\vert \tilde
{g}\left(  \mathbf{\eta},\mathbf{\lambda}\right)  \right\rangle ,
\end{equation}
where
\begin{equation}
\mathit{s}\left(  \mathbf{\eta}\right)  =\left\langle \left.  \tilde{g}\left(
\mathbf{\eta},\mathbf{\lambda}\right)  \right\vert \tilde{g}\left(
\mathbf{\eta},\mathbf{\lambda}\right)  \right\rangle =\sum_{1\leq l\leq
s}e^{2\eta_{l}}c_{0l}^{2}=\sum_{1\leq l\leq s}e^{2\eta_{l}}n_{0l}%
\end{equation}
giving
\begin{align}
\left\vert g\left(  \mathbf{\eta},\mathbf{\lambda}\right)  \right\rangle  &
=\mathit{s}\left(  \mathbf{\eta}\right)  ^{-\frac{1}{2}}\sum_{1\leq l\leq
s}e^{\eta_{l}}c_{0l}\left\vert \varphi_{l}\left(  \mathbf{\lambda}\right)
\varphi_{l+s}\left(  \mathbf{\lambda}\right)  \right\rangle \\
&  =\sum_{1\leq l\leq s}e^{\eta_{l}}c_{0l}\mathit{s}\left(  \mathbf{\eta
}\right)  ^{-\frac{1}{2}}\left\vert \varphi_{l}\left(  \mathbf{\lambda
}\right)  \varphi_{l+s}\left(  \mathbf{\lambda}\right)  \right\rangle
=\sum_{1\leq l\leq s}c_{l}\left(  \mathbf{\eta}\right)  \left\vert \varphi
_{l}\left(  \mathbf{\lambda}\right)  \varphi_{l+s}\left(  \mathbf{\lambda
}\right)  \right\rangle ,
\end{align}
where
\begin{equation}
c_{l}\left(  \mathbf{\eta}\right)  =e^{\eta_{l}}c_{0l}\mathit{s}\left(
\mathbf{\eta}\right)  ^{-\frac{1}{2}}%
\end{equation}
leading to
\begin{equation}
n_{l}\left(  \mathbf{\eta}\right)  =e^{2\eta_{l}}n_{0l}\mathit{s}\left(
\mathbf{\eta}\right)  ^{-1}.
\end{equation}
This gives
\begin{equation}
\sum_{1\leq l\leq s}n_{l}\left(  \mathbf{\eta}\right)  =\mathit{s}\left(
\mathbf{\eta}\right)  ^{-1}\sum_{1\leq l\leq s}e^{2\eta_{l}}n_{0l}=\left(
\sum_{1\leq l\leq s}e^{2\eta_{l}}n_{0l}\right)  ^{-1}\sum_{1\leq l\leq
s}e^{2\eta_{l}}n_{0l}=1
\end{equation}
which shows that $\left\{  n_{l}\left(  \mathbf{\eta}\right)  \right\}  $
satisfy the geminal occupation number positivity and normalization
constraints. We can define%
\begin{equation}
\left\vert \varphi_{l}\left(  \mathbf{\eta,\lambda}\right)  \right\rangle
=\exp\eta_{l}\exp\left\{  i\sum_{1\leq j,k\leq2s}\lambda_{jk}a_{j}^{\dagger
}a_{k}\right\}  \left\vert \varphi_{0l}\right\rangle ,
\end{equation}
so that
\begin{equation}
\left\langle \left.  \varphi_{l}\left(  \mathbf{\eta,\lambda}\right)
\right\vert \varphi_{l}\left(  \mathbf{\eta,\lambda}\right)  \right\rangle
=e^{2\eta_{l}}\left\langle \left.  \varphi_{0l}\right\vert \varphi
_{0l}\right\rangle =e^{2\eta_{l}}%
\end{equation}
and
\begin{equation}
\left\vert g\left(  \mathbf{\eta},\mathbf{\lambda}\right)  \right\rangle
=\mathit{s}\left(  \mathbf{\eta}\right)  ^{-\frac{1}{2}}\sum_{1\leq l\leq
s}c_{0l}\left\vert \varphi_{l}\left(  \mathbf{\eta,\lambda}\right)
\varphi_{l+s}\left(  \mathbf{\eta,\lambda}\right)  \right\rangle .
\end{equation}


We can express the geminal occupation numbers (and hence the occupation
numbers of the FORDM) and the normalization of the CGSO's in terms of
$\mathbf{\eta}$
\begin{align}
n_{l}\left(  \mathbf{\eta}\right)   &  =e^{2\eta_{l}}n_{0l}\mathit{s}\left(
\mathbf{\eta}\right)  ^{-1}\nonumber\\
s_{l}  &  =\left\langle \left.  \varphi_{l}\left(  \mathbf{\eta,\lambda
}\right)  \right\vert \varphi_{l}\left(  \mathbf{\eta,\lambda}\right)
\right\rangle =e^{2\eta_{l}}%
\end{align}
giving
\begin{equation}
\eta_{l}=\frac{1}{2}\ln\left\langle \left.  \varphi_{l}\left(  \mathbf{\eta
,\lambda}\right)  \right\vert \varphi_{l}\left(  \mathbf{\eta,\lambda}\right)
\right\rangle \equiv\frac{1}{2}\ln\left\langle \left.  \varphi_{l}\left(
\mathbf{\lambda}\right)  \right\vert \varphi_{l}\left(  \mathbf{\lambda
}\right)  \right\rangle =\frac{1}{2}\ln s_{l}%
\end{equation}
Hence%
\begin{equation}
\mathcal{E}\left(  \mathbf{n}_{\min}\mathbf{,\varphi}_{\min}\right)
=\mathcal{E}\left(  \mathbf{\eta}_{\min}\mathbf{,\lambda}_{\min}\right)
=\mathcal{E}\left(  \mathbf{\varphi}_{\min}\right)
\end{equation}
where
\begin{align}
\mathcal{E}\left(  \mathbf{n}_{\min}\mathbf{,\varphi}_{\min}\right)   &
=\min_{\mathbf{n\in\mathcal{S},\varphi\in}\mathcal{O}\left(  \mathbf{1}%
_{s}\right)  }\mathcal{E}\left(  \mathbf{n,\varphi}\right) \\
\mathcal{E}\left(  \mathbf{\eta}_{\min}\mathbf{,\lambda}_{\min}\right)   &
=\min_{\mathbf{\eta,\lambda}}\mathcal{E}\left(  \mathbf{\eta,\lambda}\right)
\\
\mathcal{E}\left(  \mathbf{\varphi}_{\min}\right)   &  =\min_{\mathbf{s\in
}\mathbb{R}_{+}^{s}}\min_{\mathbf{\varphi\in}\mathcal{O}\left(  \mathbf{s}%
\right)  }\mathcal{E}\left(  \mathbf{\varphi}\right)
\end{align}%
\begin{align}
\mathcal{S}  &  =\left\{  \left.  \mathbf{n}\right\vert n_{j}\geq0,1\leq j\leq
s;\sum_{1\leq j\leq s}n_{j}=1\right\} \\
\mathcal{O}\left(  \mathbf{s}\right)   &  =\left\{  \left.  \mathbf{\varphi
}\right\vert \left\langle \left.  \varphi_{j}\right\vert \varphi
_{k}\right\rangle =s_{j}\delta_{jk},1\leq j,k\leq s\right\}  .
\end{align}%
\begin{equation}
\mathbf{\lambda}=\mathbf{\lambda}^{\dagger},\mathbf{\eta\in}\mathbb{R}^{s}%
\end{equation}%
\begin{align}
\mathcal{L}\left(  \mathbf{\varphi,s}\right)   &  =\mathcal{E}\left(
\mathbf{n}_{0}\mathbf{,\varphi}\right)  -\sum_{1\leq j,k\leq2s}\varepsilon
_{kj}\left(  \left\langle \left.  \varphi_{j}\right\vert \varphi
_{k}\right\rangle -s_{j}\delta_{jk}\right) \\
\frac{\partial\mathcal{L}\left(  \mathbf{\varphi,s}\right)  }{\partial s_{j}}
&  =\varepsilon_{jj}.
\end{align}%
\begin{equation}
\mathcal{L}\left(  \mathbf{n\left(  \mathbf{\eta}\right)  ,\varphi}\right)
=\mathcal{L}\left(  \mathbf{\varphi,s}\left(  \mathbf{\eta}\right)  \right)
\end{equation}


\section{\bigskip Generalized Janaks Theorem\label{General}}

\section{\bigskip Discussion}

\appendix

\section{Reduced Density Operators \label{Denmat}}

\bigskip

\section{Operator Sets and Subspaces \label{Operators}}

\bigskip

\section{State Normalization \label{Normalization}}

If $D^{N}\in\mathcal{B}_{2}\left(  \mathcal{H}^{N}\right)  $ then
\begin{equation}
\operatorname{Tr}\left\{  \mathcal{N}D^{N}\right\}  =N\operatorname{Tr}%
\left\{  D^{N}\right\}
\end{equation}
as
\begin{equation}
\mathcal{N}\left\vert \Psi\right\rangle =N\left\vert \Psi\right\rangle
\quad\forall\left\vert \Psi\right\rangle \in\mathcal{H}^{N}%
\end{equation}
and
\begin{equation}
\operatorname{Tr}\left\{  \mathcal{N}D^{N}\right\}  =\int\gamma\left(
\mathbf{r,\xi}\right)  d\mathbf{r}d\mathbf{\xi.}%
\end{equation}
Thus
\begin{equation}
N\operatorname{Tr}\left\{  D^{N}\right\}  =N\Rightarrow\operatorname{Tr}%
\left\{  D^{N}\right\}  =1.
\end{equation}



\end{document}